% AIPL / ABCL^c+ User's Guide (English)
%   Translation of docs/aipl_guide_ja.tex (2026-08-23 edition).
%   Build:  lualatex -interaction=nonstopmode aipl_guide_en.tex   (twice)
%           (pdflatex/xelatex also work: this file has no CJK.)
\documentclass[11pt,a4paper]{article}
\usepackage{amsmath,amssymb}
\usepackage{listings}
\usepackage{xcolor}
\usepackage{geometry}
\usepackage{array}
\usepackage{booktabs}
\usepackage[hidelinks]{hyperref}
\geometry{margin=22mm}

\definecolor{codebg}{gray}{0.96}
\lstset{
  basicstyle=\ttfamily\small,
  backgroundcolor=\color{codebg},
  frame=single,
  framesep=4pt,
  rulecolor=\color{gray},
  breaklines=true,
  showstringspaces=false,
  columns=fullflexible,
  keepspaces=true,
  keywordstyle=\bfseries,
  commentstyle=\itshape\color{gray!130},
}
\lstdefinelanguage{AIPL}{
  morekeywords={class,method,var,float,new,now,future,await,reply,send,replyto,answer,
    call,if,else,while,do,select,case,timeout,self,sender,remote,print,
    true,false,int,string,bool,unit,any},
  sensitive=true,
  morecomment=[l]{//},
  morecomment=[s]{/*}{*/},
  morestring=[b]",
}
\lstset{language=AIPL}
\lstnewenvironment{term}{\lstset{language={},basicstyle=\ttfamily\footnotesize}}{}

\newcommand{\ty}[1]{\ensuremath{\mathtt{#1}}}
\newcommand{\eff}[1]{\texttt{\{#1\}}}

\title{AIPL / ABCL$^{c}$+ User's Guide\\
{\large ---\ syntax, usage, and sample programs ---}}
\author{ABCL$^{c}$+ / AIOS implementation (\texttt{yaskodama/abclcp}, branch \texttt{make-src-base})}
\date{2026-08-31 edition\\{\small Target: branch \texttt{make-src-base} (the current specification, including the seven disciplines)}}

\begin{document}
\maketitle

\begin{abstract}
This is the \textbf{current} user's guide to AIPL (ABCL$^{c}$+).
Part \ref{part:start} gives the shape of the language and the shortest path to
running something; Part \ref{part:lang} specifies the syntax and the type
system; Part \ref{part:disciplines} covers the seven disciplines that sit on
top of it; Part \ref{part:sample} walks through the sample programs with their
sources and their actual output; Part \ref{part:tool} describes the tools for
% (including running on bare-metal Raspberry Pi 3, \S\ref{sec:xinu})
building, running and checking; Part \ref{part:appendix} collects the grammar
summary, the present limits, and the migration notes from the previous edition.

AIPL is a concurrent language whose unit is the \textbf{actor}. A value is not
returned with \texttt{return} but with \texttt{reply}. Three static mechanisms
sit on top of that --- \textbf{types} (the return type is inferred from
\texttt{reply}, and can also be annotated), \textbf{effects} (a method declares
what it does to the world, as in \texttt{!\{ai, net\}}), and \textbf{deadlines}
(\texttt{now ... timeout ... else ...} bounds every wait). For example, ``a
real-time actor synchronously calls a model outside the box'' can be forbidden
by the effect check alone.

Above those sit the \textbf{seven disciplines} (Part \ref{part:disciplines}) ---
the failure type \texttt{result}$\langle\tau\rangle$, first-class reply
destinations, resource ordering, obligation levels, and session types. All of
them are built so that they \textbf{work without you writing annotations}.
As a result, \textbf{as long as the target of a wait is known as a class, or
the level of the target node is declared, deadlock freedom is guaranteed by the
type system}. That is not a claim on paper only: it is machine-checked in
Coq/Rocq.

What this guide describes is the \textbf{core} --- the part that all three
implementations (OCaml, Python, JavaScript) provide. The Python implementation
additionally has \textbf{extensions} (record types, refinement types, channels,
ownership), so the language has two layers (\S\ref{sec:migrate}).

If you only want the delta from the previous edition (2026-07-30), read
\S\ref{sec:migrate}. \texttt{become} has been \textbf{removed from the language}.
\end{abstract}

\tableofcontents

\part{Getting started}
\label{part:start}

\section{What AIPL is}

AIPL is an actor language in the lineage of ABCL/1. A program is a sequence of
\textbf{class declarations} and \textbf{global statements}; actors created from
the classes run by sending messages to each other.

\begin{lstlisting}
class Counter {
  var n = 0;                       // a field (the actor's state)
  method bump() : int !{mut} {     // return type int, effect mut
    n = n + 1;
    reply(n);                      // hand a value back to the caller
  }
}

var c = new Counter();             // create one actor
print(now c.bump() timeout 100 else 0);
\end{lstlisting}

Four differences from other languages are worth learning first.

\begin{enumerate}
\item \textbf{One actor is one thread}, and it processes the messages it
  receives one at a time. Fields are never touched concurrently.
\item \textbf{Values come back through \texttt{reply}.} There is no
  \texttt{return}. \texttt{reply} resolves the caller's future; the method
  keeps running afterwards.
\item \textbf{There are three kinds of send.} \texttt{send} (does not wait; a
  statement), \texttt{now} (waits for the reply; an expression), and
  \texttt{future} (does not wait; hands you a ticket; an expression).
\item \textbf{Waiting freezes you.} While an actor waits in a \texttt{now} it
  cannot process any other message. That is why waits carry deadlines.
\end{enumerate}

\section{Three static mechanisms}

\subsection{Types --- what you give back}

A method's return type is inferred from the \texttt{reply} calls in its body.
The optional annotation \texttt{method m(x) : T} may also be written. When you
write it, the checker verifies that every execution path replies, and the
declared type flows into expression inference as the \emph{expected} type.

\subsection{Effects --- what you do to the world}

A method declares its effects with \eff{...}. There are eight of them
(\S\ref{sec:effects}). If you omit the annotation the effects are inferred; if
you write it, ``the effects of the body $\subseteq$ the effects you declared''
is checked.

\begin{lstlisting}
method plan(goal) : string !{ai, net} { reply(ai_call("plan: " ++ goal)); }
method peek()     : int    !{}        { reply(n); }
\end{lstlisting}

The point is that \texttt{ai} and \texttt{net} are separate: on-device
inference that ``uses AI but never leaves the box'' is expressed with \eff{ai}
alone.

\subsection{Deadlines --- how long you will wait}

\begin{lstlisting}
var r = now planner.plan(g) timeout 2000 else "(timed out)";
\end{lstlisting}

If no reply arrives in time, the value of the \texttt{else} branch is used.
If you write no \texttt{else}, the expression has type
\texttt{result}$\langle\tau\rangle$ and cannot be used before it is checked
(\S\ref{sec:result}).

\subsection{And seven disciplines on top}

Types, effects and deadlines are the ground floor. Above them sit the failure
type, first-class reply destinations, resource ordering, obligation levels and
session types (Part \ref{part:disciplines}). All of them are built so that they
\textbf{work without annotations}. As a result, \textbf{as long as the target of
a wait is known as a class, or the level of the target node is declared,
deadlock freedom is guaranteed by the type system} (\S\ref{sec:levels}).

\section{Running something in five minutes}

\subsection{Building the implementation}

\begin{term}
cd ~/aios/abclcp
make -C src thread_repl
\end{term}

Building with \texttt{dune} fails because \texttt{src/lexer.ml} is generated
twice, so use \texttt{src/Makefile}.

\subsection{Running a program}

Pass a file of REPL commands with \texttt{-f}.

\begin{term}
cat > run.repl <<'EOF'
load docs/samples/guide/g2_now_future.aipl
compile
EOF
./src/abclrepl_thread -q -f run.repl < /dev/null
\end{term}

\begin{term}
twice = 43
awaited = 20
v=7
\end{term}

\textbf{Do not pass a \texttt{.aipl} file to \texttt{-f} directly.} It would be
read line by line as REPL commands, the class definitions would never be
loaded, and you would get \texttt{Actor xxx not found}. Always write
\texttt{load} and \texttt{compile}.

\subsection{Running only the type checker}

The REPL depends on SDL, so a small separate driver exists for type checking
alone (\S\ref{sec:tc}). Once built, \texttt{./tc file.aipl} prints the table of
return types, effects and signatures.

\begin{term}
$ ./tc docs/samples/guide/g6_effects.aipl
=== g6_effects.aipl ===
[OK] type check passed
[method return types (inferred from reply)]
  Counter#bump -> int
  Planner#plan -> string
  ViaNow#run -> string
[method effects (inferred / declared)]
  Counter#bump !{mut}
  Planner#plan !{ai, net}
  ViaNow#run !{ai, net}
\end{term}

\part{Language guide}
\label{part:lang}

\section{The shape of a program}

A program is a \textbf{sequence of declarations}. A declaration is one of the
following; the order is free, and names are resolved at run time.

\begin{center}
\begin{tabular}{@{}ll@{}}
\toprule
Declaration & Meaning \\
\midrule
\texttt{class C \{ ... \}}            & class definition \\
\texttt{var x = e;}                   & global variable \\
\texttt{var x = new C(args);}         & create an actor \\
\texttt{send t.m(args);}              & global asynchronous send \\
\texttt{send! t.m(args);}             & the same, with the check relaxed \\
\texttt{f(args);}                     & call a builtin function \\
\bottomrule
\end{tabular}
\end{center}

\begin{quote}\small
\texttt{if} and \texttt{while} cannot appear at global level. If you need
control flow, put it inside a method and start it with \texttt{send}.
\end{quote}

A class is a \textbf{sequence of fields} followed by a \textbf{sequence of
methods}.

\begin{lstlisting}
class Counter {
  var n = 0;                    // fields come first; declare with var or float
  float ratio = 1.5;
  method init(start) {          // init is called automatically by new
    n = start;
  }
  method bump() : int !{mut} {
    n = n + 1;
    reply(n);
  }
}
\end{lstlisting}

\begin{itemize}
\item Fields come \textbf{before} methods.
\item A class needs \textbf{at least one method} (the grammar has no
  field-only class).
\item A method named \texttt{init} is called automatically by
  \texttt{new C(args)}. If the number of arguments does not match, creation is
  skipped and \texttt{[Actor] x: no init; skipped} is printed.
\item A field's type is fixed by its initializer. \texttt{var n = 0;} is
  \ty{int}, so assigning a \ty{float} or a string to it later is a type error.
\end{itemize}

\section{Lexical structure}

\begin{center}
\small
\begin{tabular}{@{}lp{9.5cm}@{}}
\toprule
Kind & Content \\
\midrule
Comments   & \texttt{// to end of line}, \texttt{/* ... */} (not nestable) \\
Keywords   & \texttt{class}, \texttt{method}, \texttt{var}, \texttt{float},
             \texttt{new}, \texttt{send}, \texttt{send!}, \texttt{now},
             \texttt{future}, \texttt{await}, \texttt{if}, \texttt{else},
             \texttt{while}, \texttt{do}, \texttt{select}, \texttt{case},
             \texttt{timeout}, \texttt{call}, \texttt{replyto},
             \texttt{self}, \texttt{sender}, \texttt{remote},
             \texttt{true}, \texttt{false} \\
Integers   & \texttt{0}, \texttt{123} \\
Floats     & \texttt{1.5}, \texttt{100.} (a trailing dot alone makes it \ty{float}) \\
Strings    & \texttt{"..."}, with \texttt{\textbackslash\textbackslash}
             \texttt{\textbackslash"} \texttt{\textbackslash n}
             \texttt{\textbackslash t} \texttt{\textbackslash r} \\
Booleans   & \texttt{true}, \texttt{false} \\
Identifiers & \texttt{[A-Za-z\_][A-Za-z0-9\_]*} \\
\bottomrule
\end{tabular}
\end{center}

\begin{quote}\small
\texttt{float} is a keyword, so it cannot be used freely as a variable or type
name (it \emph{can} be written as a return annotation, \texttt{: float}, which
has its own grammar rule).
\end{quote}

\section{Types}

\begin{center}
\begin{tabular}{@{}ll@{}}
\toprule
Type & Description \\
\midrule
\ty{int}, \ty{float}, \ty{string}, \ty{bool}, \ty{unit} & base types \\
\ty{T[]}          & array (handled through the \texttt{array\_*} builtins) \\
\ty{future\ T}    & the result of a \texttt{future} send; \texttt{await} takes the \ty{T} out \\
\ty{actor(C)}     & an actor of class \texttt{C}; the type of \texttt{new C(...)} \\
\ty{any}          & a boundary whose content is not statically guaranteed
                    (\texttt{sender}, remote targets) \\
\bottomrule
\end{tabular}
\end{center}

A return annotation may be
\texttt{int} / \texttt{float} / \texttt{string} / \texttt{bool} /
\texttt{unit} / \texttt{any}, or a class name starting with a capital letter
(an actor type). An unknown lowercase name is a syntax error:
\texttt{unknown type name in return annotation}.

\begin{quote}\small
\textbf{Method parameters are monomorphic.} The type variables of a signature
are shared by every call site, so calling the same method with \ty{int} and
with \ty{string} conflicts. Polymorphic methods cannot be written.
\end{quote}

\begin{quote}\small
There is \textbf{no implicit conversion} between \ty{int} and \ty{float}.
If \texttt{var h = new Hello(5);} feeds a \texttt{float} field you get a type
error; write \texttt{new Hello(5.)}.
\end{quote}

\section{Expressions}

\begin{center}
\begin{tabular}{@{}ll@{}}
\toprule
Expression & Meaning \\
\midrule
\texttt{123}, \texttt{1.5}, \texttt{"abc"}, \texttt{true}, \texttt{false} & literals \\
\texttt{x} & variable, field or parameter \\
\texttt{-\,e} & unary minus; lowers to \texttt{neg(e)}. \ty{int}/\ty{float} \\
\texttt{a ++ b} & string concatenation (both sides are stringified) \\
\texttt{a + b}, \texttt{a - b}, \texttt{a * b} & arithmetic. \ty{int} with \ty{int} gives \ty{int};
  \ty{float} with \ty{float} gives \ty{float} \\
\texttt{a / b} & division. Even \ty{int} with \ty{int} gives \ty{float} (no integer division) \\
\texttt{a == b}, \texttt{a != b} & equality. \ty{int}/\ty{float}/\ty{string}/\ty{bool} \\
\texttt{a > b}, \texttt{a < b}, \texttt{a >= b}, \texttt{a <= b} & comparison; the result is \ty{bool} \\
\texttt{new C(args)} & create an actor (effect \eff{mem}) \\
\texttt{now t.m(args)} & synchronous send: wait for the reply and become that value \\
\texttt{now t.m(args) timeout $n$ else $e$} & synchronous send with a deadline \\
\texttt{future t.m(args)} & asynchronous send; returns \ty{future\ T} \\
\texttt{await e} & wait for a future to resolve \\
\texttt{await e timeout $n$ else $e'$} & \texttt{await} with a deadline \\
\texttt{f(args)} & call a builtin function \\
\texttt{( e )} & parentheses \\
\bottomrule
\end{tabular}
\end{center}

The target \texttt{t} is \textbf{a variable holding an actor reference}, or
\texttt{remote("host:port", "actorName")}. It is not a class name.

\subsection{Precedence}

From weakest to strongest:

\begin{center}
equality \texttt{==} \texttt{!=} \ $<$\ comparison \texttt{>} \texttt{<} \texttt{>=} \texttt{<=}
\ $<$\ concatenation \texttt{++} \ $<$\ additive \texttt{+} \texttt{-}
\ $<$\ multiplicative \texttt{*} \texttt{/} \ $<$\ unary minus
\end{center}

So \texttt{1 + 2 == 3} reads as \texttt{(1 + 2) == 3}, and
\texttt{"x=" ++ a + b} reads as \texttt{"x=" ++ (a + b)}.
Every binary operator is \textbf{left associative} (equality included) ---
\texttt{1 == 2 == false} reads as \texttt{((1 == 2) == false)} and evaluates to
\texttt{true}. Unary minus is a prefix, so it can be stacked
(\texttt{- - 3} is \texttt{3}).

\subsection{Do not confuse \texttt{+} with \texttt{++}}

\texttt{+} is \textbf{numeric only}. Mixing in a string gives
\texttt{no overload of + matches (string, int)}.

\begin{lstlisting}
print("count = " ++ n);      // correct
print("count = " + n);       // type error
\end{lstlisting}

\texttt{++} stringifies both sides first, so neither side has to be a string.

\subsection{Ambiguous overloads}

When both operands are unbound, \texttt{a + b} could be \ty{int} or
\ty{float}: it has no principal type. That is an error. An annotation makes the
expected type flow into the expression and settles it.

\begin{lstlisting}
method add(a, b)       { reply(a + b); }   // error: ambiguous overload
method add(a, b) : int { reply(a + b); }   // the annotation settles it
\end{lstlisting}

\texttt{AIOS\_LAX\_OVERLOAD=1} restores the old behaviour (warn and pick the
default candidate). Equality \texttt{==} and \texttt{!=} have four candidates
but all of them return \ty{bool}, so they are never ambiguous (ambiguity is
only reported when the \emph{return types} disagree).

\section{Statements}

\begin{center}
\begin{tabular}{@{}ll@{}}
\toprule
Statement & Meaning \\
\midrule
\texttt{var x = e;} & variable declaration \\
\texttt{var x = new C(args);} & declaration that creates an actor \\
\texttt{x = e;} & assignment \\
\texttt{f(args);} \quad \texttt{call f(args);} & call a builtin function \\
\texttt{send t.m(args);} & asynchronous send; does not wait for a reply \\
\texttt{send self.m(args);} & asynchronous send to yourself \\
\texttt{send sender.m(args);} & asynchronous send to the sender \\
\texttt{send! t.m(args);} & send with the check relaxed \\
\texttt{if (e) s} \quad \texttt{if (e) s else s} & conditional (the parentheses are required) \\
\texttt{while e do s} & loop (no parentheses required) \\
\texttt{\{ s ... \}} & block \\
\texttt{select \{ ... \}} & selective receive \\
\bottomrule
\end{tabular}
\end{center}

\begin{quote}\small
\texttt{now self.m()} and \texttt{now sender.m()} \textbf{cannot be written} ---
the grammar forbids them. A \texttt{now} to yourself is a self-deadlock (you are
busy running that very method, so nobody is left to produce the reply), and not
being able to write it is exactly the safe side. If you want to factor out a
computation such as inverse kinematics, put it in another actor and call
\texttt{now kin.solve(...)}.
\end{quote}

\begin{quote}\small
A variable declaration with a type annotation, \texttt{var x : T = e;}, does
\textbf{not exist yet}. It is in the design notes but unimplemented.
\end{quote}

\section{Actors and communication}

\subsection{The three sends}

\begin{center}
\begin{tabular}{@{}llll@{}}
\toprule
 & Waits? & Statement or expression & Inherits the callee's effects? \\
\midrule
\texttt{send a.m(x);}      & no          & statement & no \\
\texttt{now a.m(x)}        & until reply & expression (of the declared/inferred type) & \textbf{yes} \\
\texttt{future a.m(x)}     & no          & expression (\ty{future\ T}) & no \\
\bottomrule
\end{tabular}
\end{center}

``Only the side that waits carries the callee's effects'' is the rule of effect
propagation (\S\ref{sec:effects}).

\subsection{\texttt{reply}}

A method does not ``return'' a value; \texttt{reply(v)} resolves the caller's
future. The type system gives each method a return type, and every
\texttt{reply} in the body together with every \texttt{now} / \texttt{future}
call site shares that type.

\begin{itemize}
\item \texttt{reply} happens \textbf{at most once}. A second one would arrive
  after the waiter already took the first value, so it would be silently
  dropped. That is made a type error.
\item If the return type is \emph{declared} as something other than \ty{unit},
  the checker verifies that \textbf{every execution path replies exactly once}.
  The body of a \texttt{while} may run zero times, so it is conservatively
  judged as ``there is a path that does not reply''.
\item A method with no \texttt{reply} at all is treated as returning \ty{unit}.
  Notification-style methods called with \texttt{send} are of this kind.
\end{itemize}

\begin{lstlisting}
method m(k) : int { reply(k); reply(k + 1); }                     // type error
method m(k) : int { if (k > 0) { reply(1); } else { reply(2); } } // OK
\end{lstlisting}

\subsection{\texttt{select}}

\begin{lstlisting}
select {
  case job(k) -> { reply("done:" ++ k); }
  timeout 500 -> { print("timed out"); }
}
\end{lstlisting}

\texttt{case m(x1, ..., xn)} is matched against the signature of the method of
the same name. A mismatch in the number or type of arguments is a type error.
\texttt{timeout} is in milliseconds and may be omitted.

\textbf{A \texttt{reply} in a case body answers ``the selected message'', not
the enclosing method.} The implementation switches the reply destination to the
selected message before running the case, so the \texttt{reply} above answers
\texttt{job}. Hence a method containing a \texttt{select} may itself be
\ty{unit}. The \texttt{timeout} body, by contrast, runs with the enclosing
method's reply destination.

\begin{quote}\small
\texttt{select} looks for messages that have \emph{not yet been processed}. If
ordinary method dispatch takes the message of that name first, nothing is left
for the case, so it is safer to design the names you \texttt{select} on so that
nobody calls them directly (see the output in \S\ref{sec:g4}).
\end{quote}

\section{Waiting with a deadline}
\label{sec:deadline}

\subsection{Why it is needed}

A \texttt{now} or \texttt{await} without a deadline waits forever if no reply
comes. An actor is one thread and a wait blocks on a condition variable, so
while it waits it cannot process any other message. It turns to stone.

\begin{lstlisting}
now t.m(a) timeout <milliseconds> else <expression>
await f    timeout <milliseconds> else <expression>
\end{lstlisting}

If the reply arrives in time you get its value; otherwise you get the value of
the \texttt{else} branch.

\subsection{The rules}

\begin{itemize}
\item The \texttt{else} branch must have \textbf{the same type as the body},
  because a timeout must not change the type of the whole expression.
  Otherwise: \texttt{now: the else branch has type string but the call returns int}.
\item The deadline must be \textbf{positive} (\texttt{timeout must be positive}).
\item A \texttt{now} / \texttt{await} without a deadline is a \textbf{warning}
  by default, and an \textbf{error} under \texttt{AIOS\_STRICT\_DEADLINE=1}.
\item The upper bound on the number of replies is the \texttt{max} of the body
  and the \texttt{else} branch (not their sum), since only \emph{one} of them
  runs.
\item A rejected future (\texttt{reject}) is distinguished from a timeout and
  raises an exception.
\end{itemize}

In the implementation, OCaml's \texttt{Condition} has no timed wait, so the
deadline is measured by polling at 1\,ms granularity.

\subsection{Relation to type soundness}
\label{sec:soundness}

Writing a deadline makes every wait finish in finite time. But that only means
\emph{a deadlock is converted into a failure in finite time}; the failure does
not disappear. Two actors waiting on each other with \texttt{now} both time
out. Nothing is stuck --- and nothing is right either.

To say \textbf{that nothing gets stuck} in the type system you need not
deadlines but \textbf{obligation levels} (\S\ref{sec:levels}). Those are
\textbf{machine-checked in Coq/Rocq}: for the core \emph{including}
\texttt{await}, type soundness, type safety and deadlock freedom are proved
with no axioms (\texttt{coq/AIPLSoundness2.v}). The first version could state
deadlock freedom only for a fragment with \texttt{await} removed, which is what
the second version fixes.

The progress theorem changed as follows.

\begin{center}
\begin{tabular}{@{}ll@{}}
\toprule
& What progress states \\
\midrule
1st version & terminal state $\lor$ can step $\lor$ \textbf{all tasks are waiting} \\
2nd version & terminal state $\lor$ can step \\
\bottomrule
\end{tabular}
\end{center}

Termination is still not claimed. \texttt{while 1 > 0 do \{\}} does not stop,
but it is not a wait, so it does not contradict progress.

\section{Effects}
\label{sec:effects}

\subsection{The eight effects}

\begin{center}
\begin{tabular}{@{}llll@{}}
\toprule
Effect & Meaning & Effect & Meaning \\
\midrule
\texttt{mut}  & assignment to a field  & \texttt{net} & communication outside the node \\
\texttt{time} & reading the clock, waiting & \texttt{ai}  & model inference \\
\texttt{io}   & input/output to a device & \texttt{fs}  & persistence \\
\texttt{mem}  & dynamic allocation     & \texttt{log} & output for observation only \\
\bottomrule
\end{tabular}
\end{center}

This is not a newly designed taxonomy: it is the capability classification the
evaluator already carried as run-time metadata, made static. The point is that
\texttt{ai} and \texttt{net} are separate --- a model call that leaves the box
has both, while on-device inference gets \eff{ai} only. ``Uses AI but never
leaves the box'' becomes visible in the signature.

\subsection{How to write them}

\begin{lstlisting}
method plan(goal) : string !{ai, net} { ... }   // after the return type
method peek()            !{}          { ... }   // explicitly effect-free
method tick()                         { ... }   // omitted: inferred
\end{lstlisting}

The annotation is optional (gradual). If you omit it, inference simply runs and
nothing is said. If you write it, ``the effects of the body $\subseteq$ the
effects declared'' is checked. Declaring more than you use is allowed. An
unknown effect name is rejected as a typo.

\begin{term}
[Type error] line 1, col 18: unknown effect(s) network in A.m;
  known effects are mut, time, io, mem, net, ai, fs, log
\end{term}

\subsection{What counts as an effect}

\begin{center}
\small
\begin{tabular}{@{}ll@{}}
\toprule
Construct & Effect \\
\midrule
assignment to a field        & \eff{mut} \\
assignment to a local        & \textbf{not an effect} \\
\texttt{new C(...)}          & \eff{mem} \\
calling a builtin            & that function's effects (\S\ref{sec:builtins}) \\
send to a \texttt{remote(...)} target & \eff{net} \\
local send with \texttt{now} & \textbf{inherits the callee's effects} \\
\texttt{send} / \texttt{future} & does not inherit \\
\bottomrule
\end{tabular}
\end{center}

That \texttt{mut} covers only assignment to fields is deliberate: it lets a
pure actor (safe to replicate) and an actor with external effects (exactly one
at a time) be told apart by the effect set alone.

Propagation through \texttt{now} iterates to a fixpoint. Effects only grow, so
it terminates.

\subsection{What this lets you forbid}

\begin{lstlisting}
class Controller {
  method step(s) : string !{io, time} {
    var r = now planner.plan(s) timeout 100 else "";   // plan is !{ai, net}
    reply(r);
  }
}
\end{lstlisting}

\begin{term}
[Type error] method Controller.step declares {io, time}
  but its body has {ai, net, io, time} (missing {ai, net})
\end{term}

``A real-time actor synchronously calls an agent outside the box'' is forbidden
by the effect check alone. What used to be stated as a design decision now
holds mechanically in the type system.

\begin{quote}\small
\textbf{A known hole.} Splitting a \texttt{now} into \texttt{future} and
\texttt{await} \emph{escapes} this check: the implementation adds the
propagation edge only at the \texttt{now} site and not at the \texttt{await}.
A reproduction is in
\texttt{docs/samples/}\allowbreak\texttt{soundness2\_effect\_gap.aipl}.
The fix (carrying effects in the \ty{future} type) is described in the second
soundness report.
\end{quote}

\part{The seven disciplines}
\label{part:disciplines}

On top of types, effects and deadlines sit seven disciplines. All of them are
built so that they \textbf{work without annotations} --- because the
information they need is usually already written in the program.

\begin{center}
\small
\begin{tabular}{@{}lp{5.6cm}p{4.0cm}@{}}
\toprule
Discipline & How you write it & What it stops \\
\midrule
Effects & \texttt{!\{ai, net\}} & the real-time side waiting on a model outside the box \\
Deadlines & \texttt{timeout $n$ else $e$} & waits that never return \\
Failure type & \texttt{timeout $n$} (no \texttt{else}) & letting a timed-out value flow on as a normal one \\
First-class reply & \texttt{replyto} / \texttt{answer} / parameter type \texttt{reply} & dropped replies and double replies \\
Resource order & \texttt{acquire} / \texttt{release} / \texttt{resource\_order} & forgetting to release, double acquire, acquiring in reverse order \\
Obligation levels & \texttt{@$n$} (inferred if omitted) & cycles of waiting \\
Session types & \texttt{protocol\_define} / \texttt{\_start} / \texttt{\_end} & steps taken out of order \\
\bottomrule
\end{tabular}
\end{center}

\section{Putting failure in the type --- \texttt{result}$\langle\tau\rangle$}
\label{sec:result}

If you write \texttt{else} on a deadline, the timed-out value and the normal
value have \textbf{the same type}. The program cannot ask whether $-1$ came
back because it timed out or because $-1$ was computed.

Write no \texttt{else} and you get \texttt{result}$\langle\tau\rangle$.

\begin{center}
\begin{tabular}{@{}ll@{}}
\toprule
How you write it & Type \\
\midrule
\texttt{now a.m(x) timeout 200 else -1} & \ty{int} \\
\texttt{now a.m(x) timeout 200} & \texttt{result}$\langle$\ty{int}$\rangle$ \\
\bottomrule
\end{tabular}
\end{center}

Use \texttt{is\_ok} / \texttt{timed\_out} / \texttt{value} to take the value out.

\begin{lstlisting}
var a = now s.fast(42) timeout 200;      // result<int>
if (is_ok(a)) { print("fast ok " ++ value(a, -1)); }
else          { print("fast timed out"); }
\end{lstlisting}

A \texttt{result}$\langle\tau\rangle$ cannot be used as a number. Feeding it
into arithmetic without checking stops the program at compile time.

\begin{term}
no overload of + matches (result int, int)
\end{term}

\section{Making the reply destination a first-class value}
\label{sec:reply1st}

\texttt{replyto} yields, as a value, the reply destination of the message
currently being processed. Its type is \texttt{reply}$\langle\rho\rangle$
(where $\rho$ is the method's return type), and \texttt{answer(r, v)} sends the
reply. \texttt{reply(v)} is sugar for that.

With the parameter annotation \texttt{reply} you can hand it to another actor.
A front desk takes the request and a specialist answers \textbf{directly}.

\begin{lstlisting}
class Backend {
  method handle(r: reply, x) : unit !{} {
    answer(r, x * 2);            // answer Frontend's caller directly
  }
}
class Frontend {
  var b = new Backend();
  method ask(x) : int !{} {
    send b.handle(replyto, x);   // pass the destination on; do not reply here
  }
}
\end{lstlisting}

A reply destination is \textbf{linear}: using it exactly once is an obligation.

\begin{center}
\small
\begin{tabular}{@{}lp{7.6cm}@{}}
\toprule
Rule & Content \\
\midrule
Creation & \texttt{var r = replyto;} creates the obligation \\
Consumption & \texttt{answer(r, v)}, or passing it as a send argument, discharges it \\
Exactly once & the same \texttt{r} must not be used twice \\
Nothing left & no obligation may remain when the method returns \\
Branches & the two branches of an \texttt{if} must join in the same state \\
\bottomrule
\end{tabular}
\end{center}

\begin{quote}\small
The check is closed within a method. A reply destination \textbf{cannot be kept
in a field} (a field's type is fixed by its initializer). A bounded buffer that
``collects senders and answers them later'' therefore cannot be written
directly in the core as it stands.
\end{quote}

\section{Resource order}
\label{sec:resorder}

Effects are a \textbf{set}, so they cannot express order; ``take it and give it
back'' cannot be said. So the \texttt{acquire} / \texttt{release} pairs are
tracked through the body instead.

\begin{lstlisting}
class Bus {
  var n = 0;
  method read(addr) : int !{mut} {
    acquire("i2c");
    n = n + 1;
    var v = addr * 2;
    release("i2c");           // delete this line and it is a type error
    reply(v);
  }
}
\end{lstlisting}

There are three rules. Nothing may still be held when the method returns. You
may not release what you do not hold. You may not acquire twice. The two
branches of an \texttt{if} must join holding the same things, and the body of a
\texttt{while} must not change what is held.

\subsection{A global order --- stopping reverse acquisition}

Checking pairs is not enough. If two actors take the same two resources in
\textbf{opposite order}, both pairings are correct, yet at run time they wait
for each other.

No new annotation is needed. The nesting ``$s$ was acquired while $r$ was
held'' \emph{is} the claim $r < s$. The claims are collected from the whole
program and checked for cycles.

\begin{term}
resource order: cache -> db -> cache is circular;
  cache before db (in Reader.run); db before cache (in Writer.run).
  Acquiring in opposite orders deadlocks
\end{term}

Every edge of the cycle carries \textbf{where it came from}, so the two places
that acquire in opposite order are named directly.

You may also write \texttt{resource\_order("bus -> dev")} instead of relying on
inference. A declaration simply enters the same table as an edge, so if the
declaration and the actual acquisitions disagree, that shows up as a cycle too.
The inferred order can be printed with \texttt{AIOS\_SHOW\_LEVELS=1}.

\begin{term}
[resource] bus                  @0
[resource] dev                  @1
\end{term}

\begin{quote}\small
Only \texttt{acquire} calls whose resource name is a \textbf{string literal}
can be tracked. A name held in a variable is invisible statically, so the
declared order is also enforced at run time
(\texttt{AIOS\_STRICT\_RESOURCE=1} reports the places that could not be tracked).
\end{quote}

\section{Obligation levels --- stopping cycles of waiting}
\label{sec:levels}

There are four shapes in which a wait never returns: a cycle of waits, a wait
that is never replied to, a \texttt{select} waiting for a message that never
comes, and resources acquired in reverse order. \textbf{Obligation levels}
address the first.

Each method gets a number, and \textbf{a wait must always go to a strictly
larger number}.

\begin{lstlisting}
class Leaf { method f(x)  : int @2 !{} { reply(x + 1); } }
class Mid  { var l = new Leaf();
             method g(x)  : int @1 !{} { reply(now l.f(x) timeout 200 else 0); } }
\end{lstlisting}

\textbf{You do not have to write them.} If you omit them, they are inferred by
a fixpoint that pushes levels up along the wait edges. An explicit annotation
is treated as a fixed point; if it cannot be satisfied, that is where the
contradiction is.

\begin{term}
obligation level: B9#g (@2) waits on A9#f (@2);
                  a wait must go to a strictly higher level
obligation levels do not converge; the wait graph has a cycle
\end{term}

The inferred levels can be printed with \texttt{AIOS\_SHOW\_LEVELS=1}.

\subsection{Waits that cross nodes}

In a mesh the target lives on another node and is invisible statically. So each
node declares a \textbf{floor} for its levels, and a wait towards that node
must go up.

\begin{lstlisting}
node_level("rt0", 10);            // the real-time node is a leaf (high level)
\end{lstlisting}

\begin{term}
obligation level: Ctl2#go (@7) waits on node rt0 (floor @1);
                  a wait across nodes must go up
\end{term}

This is exactly the same shape as checking effects at the border with
\texttt{node\_allow}.

\subsection{How much is guaranteed}

\textbf{As long as the target of a wait is known as a class, or the level of the
target node is declared, deadlock freedom is guaranteed by the type system.}
This is not a claim on paper only --- it is machine-checked in Coq/Rocq for the
core \emph{including} \texttt{await} (\texttt{coq/AIPLSoundness2.v}, no axioms).

The conditions matter. A dynamic target is invisible. Levels are per method, so
two instances of the same class are treated alike. Termination and starvation
are separate problems.

\section{Session types --- the order of an exchange}
\label{sec:session}

``Open it, then read, then close it at the end'' is a promise that neither
deadlines nor the existence of a sender can express. AIPL already had
\textbf{run-time} session protocols, so the type checker reads
\textbf{the same declaration}. No new declaration is needed.

\begin{lstlisting}
protocol_define("Pipeline", "f.get -> r.scale -> st.put");
var sid = protocol_start("Pipeline");

var a = now f.get(4)    timeout 200 else 0;
var b = now r.scale(a)  timeout 200 else 0;
var c = now st.put(b)   timeout 200 else 0;
protocol_end(sid);
\end{lstlisting}

Get the order wrong and it is reported before the program runs.

\begin{term}
protocol Pipeline: expected r.scale but the program sends st.put here
protocol Pipeline is incomplete at protocol_end; next expected st.put
\end{term}

(An example of real output.)

\begin{term}
protocol P29: expected b.step2 but the program sends c.step3 here
\end{term}

\subsection{When the session crosses actors}

In a real program the steps do not fit in one straight sequence: you call one
front-desk actor, and the remaining steps happen inside its body.

A synchronous send (\texttt{now} / \texttt{aios\_now}) \textbf{runs to
completion before} the caller continues. So the callee's sequence of sends can
be spliced in at the call site (behaviour expansion), and a session that
crosses actors can still be matched as a single sequence.

Asynchronous sends (\texttt{send} / \texttt{future}) are not spliced in ---
there is no telling when they run. If the callee contains steps, completion is
left to the run-time check (\texttt{AIOS\_STRICT\_PROTOCOL=1} reports it).

\begin{quote}\small
Steps refer to roles \textbf{by name}, so \textbf{pass actors as arguments}.
If you keep them in fields, the run-time name (\texttt{c\#f}) and the static
name (\texttt{f}) disagree and the step cannot point at the role.
\end{quote}

\section{The \texttt{select} discipline}
\label{sec:selectrule}

Two things are required of a \texttt{select}.

\begin{enumerate}
\item \textbf{Write a deadline.} A \texttt{select} without a \texttt{timeout}
  clause waits forever for a message that may never come.
\item \textbf{Check that a sender exists.} If nobody in this program sends the
  \texttt{m} of \texttt{case m(...)}, it is either a typo or something you
  forgot to write.
\end{enumerate}

\begin{term}
select waits for W2.ghost but nothing in this program sends it
select without a timeout waits forever; write `timeout <ms> -> { ... }`
\end{term}

Both are warnings by default. When sends are dynamic, or the program is exposed
to the outside, no judgement is possible, so the check is simply not made ---
\textbf{the language does not see the whole world}.

\section{Mesh deployment}
\label{sec:mesh}

An actor's source is shipped to another node and is parsed and type-checked
\textbf{at the destination} before it is instantiated. The effect annotations
travel with the code and are matched against the receiving node's policy.

\begin{lstlisting}
node_allow("rt0", "mut, log, time");     // effects the real-time node accepts
deploy("rt0", "Servo", "servo1");        // ship Servo's source to rt0 and instantiate
var d = now remote("rt0", "servo1").step(3) timeout 200 else -1;
\end{lstlisting}

\begin{term}
[deploy] Servo -> rt0/servo1 (compiled at destination)
\end{term}

Ship code carrying \texttt{ai} or \texttt{net} to a real-time node and it is
stopped at the border. Effects turn from a discipline inside one node into
\textbf{immigration control between nodes}.

\section{Escape hatches (environment variables)}
\label{sec:escape}

Every default is on the strict side. Environment variables exist that loosen
things, but they are not meant for everyday use. The list is in
\S\ref{sec:envvars}.

\part{Reference}

\section{What the type checker looks at}

\begin{center}
\small
\begin{tabular}{@{}lp{9.0cm}@{}}
\toprule
Check & Content \\
\midrule
Method existence   & does the local target have a method of that name \\
Argument count and type & do the actual arguments match what the body requires \\
\texttt{reply} type & do all \texttt{reply} calls have the same (declared) type \\
\texttt{reply} count & at most once; with an annotation, exactly once on every path \\
select matching    & do the \texttt{case} arity and argument types match the signature \\
Annotations when exposed & do the methods of a \texttt{web\_expose}d actor carry annotations \\
Overloads          & is the candidate unique (ambiguity is an error) \\
Effects            & do the declared effects cover the body's; are the effect names known \\
Deadlines          & the type of the \texttt{else} branch, positivity, presence \\
Failure type       & is a \texttt{result}$\langle\tau\rangle$ used without being checked \\
Reply linearity    & is \texttt{replyto} used exactly once \\
\texttt{reply} totality & does a method waited on by \texttt{now}/\texttt{await} reply on every path \\
Resources          & the \texttt{acquire}/\texttt{release} pairs and cycles in the acquisition order \\
Obligation levels  & does every wait go up; the same across node borders \\
Sessions           & is the order of sends the order of \texttt{protocol\_define} \\
select             & is there a deadline; does anybody send the awaited message \\
Assignment to undeclared names & is a name assigned to without a \texttt{var} declaration \\
Deployment border  & do the shipped code's effects fit the receiver's policy \\
\bottomrule
\end{tabular}
\end{center}

Representative messages:

\begin{term}
no method foo in actor(C)
type mismatch                                  <- argument or assignment type
no overload of + matches (string, int)         <- + is numeric only
ambiguous overload of + for ('a, 'a): candidate return types are float / int
reply type mismatch: this method is declared to reply with int,
  but replies with string here
method m is declared to reply with int, but some execution path does not reply
method m may reply more than once on some path
select: case m binds 1 variable(s) but method m takes 2
actor C is reachable from outside this program, but m has no return type
  annotation; ...
method C.m declares {log, time} but its body has {ai, log, net}
  (missing {ai, net})
unknown effect(s) network in A.m; known effects are mut, time, io, mem, ...
now: the else branch has type string but the call returns int
now: timeout must be positive
now without a deadline waits forever;
  write `now ... timeout <ms> else <expr>`      <- warning by default
assignment to undeclared name: cout (write `var cout = ...` to declare it)
no overload of + matches (result int, int)     <- check a result before using it
method Silent4.f is waited on by now/await but does not reply on every path
resource i2c is released without being acquired
method Dev4.use returns while still holding i2c
resource order: cache -> db -> cache is circular; ...
obligation level: B9#g (@2) waits on A9#f (@2);
  a wait must go to a strictly higher level
obligation levels do not converge; the wait graph has a cycle
protocol P29: expected b.step2 but the program sends c.step3 here
select waits for W2.ghost but nothing in this program sends it   <- warning
select without a timeout waits forever                            <- warning
deploy: Servo.step is @3 but node rt0 has floor @10
\end{term}

\section{Exposing to the outside}
\label{sec:expose}

\texttt{web\_listen(port)} starts an HTTP gateway, and
\texttt{web\_expose(path, actorName)} gives an exposed endpoint an alias.

\begin{center}
\begin{tabular}{@{}ll@{}}
\toprule
Endpoint & Content \\
\midrule
\texttt{GET /}                & a small UI \\
\texttt{POST /api/send}       & \texttt{to=<actor>\&method=<m>\&args=<a,b>} reaches any actor \\
\texttt{POST /api/x/<key>}    & reaches the alias registered with \texttt{web\_expose} \\
\bottomrule
\end{tabular}
\end{center}

\texttt{POST /api/send} \textbf{can reach any actor by name}. In other words it
is \texttt{web\_listen} that opens the border; \texttt{web\_expose} only adds an
alias. The type checker takes that into account: for an actor named by
\texttt{web\_expose}, a return annotation on every method (except \texttt{init})
is \textbf{mandatory}, and with \texttt{web\_listen} alone it is a
\textbf{warning}. The far side cannot see the body, so the reply type can only
travel as a declaration.

\texttt{AIOS\_LAX\_EXPOSE=1} downgrades the error to a warning.

Arguments arriving over HTTP are built by interpreting the string as \ty{int}
$\to$ \ty{float} $\to$ \ty{string}, in that order. An annotation cannot make
the boundary less dynamic, but it can declare \emph{which type you promise}.

\section{Builtin functions and their effects}
\label{sec:builtins}

\begin{center}
\small
\begin{tabular}{@{}lp{9.4cm}l@{}}
\toprule
Group & Names & Effect \\
\midrule
Math & \texttt{sin}, \texttt{cos}, \texttt{tan}, \texttt{asin}, \texttt{acos},
  \texttt{atan}, \texttt{sqrt}, \texttt{exp}, \texttt{log10}, \texttt{abs},
  \texttt{floor}, \texttt{ceil}, \texttt{round} & \eff{} \\
Basic & \texttt{typeof}, \texttt{reply}, \texttt{neg} & \eff{} \\
Output & \texttt{print} & \eff{log} \\
Time & \texttt{wait} & \eff{time} \\
Arrays (read) & \texttt{array\_get}, \texttt{array\_len} & \eff{} \\
Arrays (allocate) & \texttt{array\_empty}, \texttt{array\_push}, \texttt{array\_set} & \eff{mem} \\
Actor creation & \texttt{spawn} & \eff{mem} \\
Drawing (SDL) & \texttt{sdl\_init}, \texttt{sdl\_clear}, \texttt{sdl\_line},
  \texttt{sdl\_erase\_line}, \texttt{sdl\_present}, \texttt{sdl\_line\_c} & \eff{io} \\
Memory and tasks & \texttt{aios\_memory\_put/get/has/keys},
  \texttt{aios\_task\_create/set/get/info}, \texttt{aios\_tasks} & \eff{fs} \\
Model inference & \texttt{ai\_call}, \texttt{ai\_call\_with\_system},
  \texttt{model\_generate}, \texttt{gemini\_generate},
  \texttt{openai\_generate} & \eff{ai, net} \\
Communication and exposure & \texttt{web\_listen}, \texttt{web\_expose},
  \texttt{aios\_now}, \texttt{aios\_future},
  \texttt{remote\_review}, \texttt{remote\_review\_ja},
  \texttt{remote\_reviewer\_host} & \eff{net} \\
Introspection & \texttt{capabilities}, \texttt{capability\_prims},
  \texttt{aios\_kernel}, \texttt{aios\_actors}, \texttt{aios\_actor\_info},
  \texttt{aios\_actor\_methods}, \texttt{aios\_mailbox\_len},
  \texttt{actor\_dump} & \eff{log} \\
Events & \texttt{aios\_emit}, \texttt{aios\_events},
  \texttt{aios\_events\_since}, \texttt{aios\_event\_count} & \eff{log} \\
Protocols & \texttt{protocol\_define/start/use/current/state/end/events} & \eff{log} \\
Services & \texttt{aios\_register\_service}, \texttt{aios\_services},
  \texttt{aios\_service\_actor}, \texttt{aios\_service\_info} & \eff{log} \\
\bottomrule
\end{tabular}
\end{center}

A primitive that is not in this table is treated as effect-free.

\part{Sample programs}
\label{part:sample}

The samples in this part are the ten programs in \texttt{docs/samples/guide/}.
All of them type-check, their output is exactly as printed here, and
\textbf{the three implementations (OCaml, Python, JavaScript) agree on that
output} --- verified with \texttt{tools/run\_all\_impls.sh}.
(The sources in the repository carry Japanese comments; the listings below are
the same programs with the comments translated.)

Samples for the more advanced features (\texttt{result}$\langle\tau\rangle$,
delegating a reply destination, resource order, obligation levels, session
types, mesh deployment) are the twenty programs in
\texttt{docs/samples/paper/}; the examples in Part \ref{part:disciplines} come
from there.

\section{g1: classes, fields, \texttt{send}, string concatenation}

\begin{lstlisting}
// g1: class, field, init, method, send, string concatenation
class Greeter {
  var count = 0;              // a field, initialized as int

  method init(name) {         // called automatically by new
    print("hello, " ++ name); // ++ concatenates (both sides are stringified)
  }

  method tick() : unit {      // return annotation; no reply, so unit
    count = count + 1;        // + is numeric only
    print("tick " ++ count);
  }
}

var g = new Greeter("AIPL");
send g.tick();                // asynchronous; does not wait for a reply
send g.tick();
\end{lstlisting}

\begin{term}
hello, AIPL
tick 1
tick 2
\end{term}

\texttt{init} is called automatically by \texttt{new Greeter("AIPL")}.
\texttt{tick} does not \texttt{reply}, so it declares \ty{unit} and the
coverage check does not apply. \texttt{send} does not wait for a reply, so the
two \texttt{tick}s are processed in order. \texttt{count + 1} is \ty{int} with
\ty{int}, so it stays \ty{int} and prints as \texttt{1}, \texttt{2}.

The inferred effects are \texttt{Greeter\#init !\{log\}} and
\texttt{Greeter\#tick !\{log, mut\}}.

\section{g2: \texttt{now} / \texttt{future} / \texttt{await} and annotations}

\begin{lstlisting}
// g2: now / future / await and the return annotation
class Calc {
  // with an annotation, every reply is matched against this type and
  // replying on every path is checked as well
  method twice(a) : int {
    reply(a * 2);
  }

  // without one, the type is inferred from reply (-> string)
  method label(a) {
    reply("v=" ++ a);
  }
}

var c = new Calc();

var s = now c.twice(21);      // now is an expression; the result is int
print("twice = " ++ (s + 1)); // usable as int in arithmetic

var f = future c.twice(10);   // future gives future int
var v = await f;              // await takes the value out
print("awaited = " ++ v);

print(now c.label(7));        // the inferred string
\end{lstlisting}

\begin{term}
twice = 43
awaited = 20
v=7
\end{term}

\texttt{twice} declares \ty{int}, so \texttt{now c.twice(21)} is \ty{int} and
can be used in the arithmetic \texttt{s + 1}. \texttt{label} omits the
annotation, and \ty{string} is inferred from \texttt{reply("v=" ++ a)}.

This sample writes no deadlines, so the type checker emits three warnings.

\begin{term}
[warn] line 16, col 9: now without a deadline waits forever;
  write `now ... timeout <ms> else <expr>`
\end{term}

\section{g3: several actors, and a \texttt{now} across actors}

\begin{lstlisting}
// g3: several actors, and a now that crosses actors
class Stock {
  var n = 10;

  method take(k) : int {
    n = n - k;
    if (n < 0) { n = 0; }
    reply(n);
  }
}

class Front {
  // Stock.take returns int; that becomes the type of the now,
  // which in turn feeds place's own string
  method place(k) : string {
    var left = now stock.take(k);
    if (left > 0) {
      reply("ok, left=" ++ left);
    } else {
      reply("sold out");
    }
  }
}

var stock = new Stock();
var front = new Front();

print(now front.place(3));
print(now front.place(99));
\end{lstlisting}

\begin{term}
ok, left=7
sold out
\end{term}

\texttt{Stock.take} returns \ty{int}; that becomes the type of the \texttt{now}
inside \texttt{Front.place}, which then feeds \texttt{place}'s own \ty{string}.
Types propagate across actor boundaries.

Because \texttt{place} declares \ty{string}, the checker verifies that both
branches of the \texttt{if} reply. Delete one and you get
\texttt{some execution path does not reply}.

Note that \texttt{n = n - k} with \ty{int} \texttt{n} forces \texttt{k} to be
\ty{int}, consistent with \texttt{place(3)} at the call site; writing
\texttt{place("x")} would be an \emph{argument} type error.

Both methods have effect \eff{mut} (assignment to field \texttt{n}), and
\texttt{Front\#place} inherits \texttt{Stock\#take}'s \eff{mut} through the
\texttt{now}.

\section{g4: \texttt{select}, and which message a \texttt{reply} answers}
\label{sec:g4}

\begin{lstlisting}
// g4: select / case / timeout, and which message a case's reply answers
class Worker {
  // the job message that arrives from outside
  method job(k) : string {
    reply("direct:" ++ k);
  }

  // wait for job with select. The reply in the case body answers job,
  // so serve itself replies to nothing and stays unit
  method serve() : unit {
    print("waiting");
    select {
      case job(k) -> {
        print("got " ++ k);
        reply("via select:" ++ k);
      }
      timeout 500 -> {
        print("timed out");
      }
    }
  }
}

var w = new Worker();
send w.serve();
print(now w.job(1));
\end{lstlisting}

\begin{term}
direct:1
waiting
timed out
\end{term}

The \texttt{reply} in the body of \texttt{case job(k)} answers \texttt{job},
not \texttt{serve}, so \texttt{serve} may stay \ty{unit}. The type checker also
matches that \texttt{reply} against \texttt{job}'s return type.

The output shows \textbf{a design caveat this example carries}. Ordinary method
dispatch processed \texttt{now w.job(1)} first, so by the time \texttt{serve}'s
\texttt{select} ran, no \texttt{job} was left in the mailbox and it fell
through to the \texttt{timeout}. It is safer to design the message names you
\texttt{select} on so that they are not called directly.

\section{g5: exposing to the outside makes annotations mandatory}

\begin{lstlisting}
// g5: exposing an actor makes the return annotation mandatory
//     web_listen opens the border, and POST /api/send can reach any actor.
//     The far side cannot see the body, so the reply type can only
//     travel as a declaration.
class Echo {
  method say(msg) : string {   // without the annotation this is a type error
    reply("echo: " ++ msg);
  }
  method note(msg) : unit {
    print("[note] " ++ msg);
  }
}

var echo = new Echo();

web_listen(8080);
web_expose("/echo", "echo");

print("try POST /api/x/echo with method=say&args=hi");
\end{lstlisting}

This program type-checks. Remove the annotation on \texttt{say} and you get:

\begin{term}
[Type error] line 16, col 1: actor Echo is reachable from outside this
  program, but say has no return type annotation; a remote caller cannot
  see the body, so the reply type has to be declared (method say(...) : T)
\end{term}

Running \texttt{web\_listen(8080)} starts the gateway and blocks listening, so
check it from another terminal.

\begin{term}
$ curl -s -X POST http://localhost:8080/api/x/echo \
       -d 'method=say&args=hi'
\end{term}

\section{g6: effect annotations and propagation}

\begin{lstlisting}
// g6: the effect annotation !{...}
//
//     The existing capability classification, made static. Eight effects:
//       mut  write your own state        net  communication outside the node
//       time reading the clock, waiting  ai   model inference
//       io   input/output to a device    fs   persistence
//       mem  dynamic allocation          log  output for observation only
//
//     Omit the annotation and it is inferred from the body. Write it and
//     "the body's effects  the declared effects" is checked (more is allowed).
class Counter {
  var n = 0;

  method bump() : int !{mut} {   // assignment to a field is mut
    n = n + 1;
    reply(n);
  }
  method peek() : int !{} {      // effect-free: pure, so free to replicate
    reply(n);
  }
  method show() : unit !{log} {  // print is log
    print("n = " ++ n);
  }
}

class Planner {
  // ai_call leaves the box, so it has both ai and net.
  // It shows up in the signature, so which actors leave the box is visible
  method plan(goal) : string !{ai, net} {
    reply(ai_call("plan: " ++ goal));
  }
}

class ViaNow {
  // now waits, so it inherits the callee's effects.
  // Without declaring {ai, net} this is a type error
  method run(g) : string !{ai, net} {
    var r = now planner.plan(g);
    reply(r);
  }
}

class ViaSend {
  // send does not wait, so it does not inherit. Effect-free is fine
  method fire(g) : unit !{} {
    send planner.plan(g);
  }
}

var planner = new Planner();
var c = new Counter();
var a = new ViaNow();
var b = new ViaSend();

// ai_call would hit an external API, so only Counter is exercised here
print(now c.bump());
print(now c.peek());
send c.show();
\end{lstlisting}

\begin{term}
1
1
n = 1
\end{term}

The effect table printed by the type checker is the point of this sample.

\begin{term}
[method effects (inferred / declared)]
  Counter#bump !{mut}          <- assignment to a field
  Counter#peek !{}             <- pure
  Counter#show !{log}          <- print
  Planner#plan !{ai, net}      <- ai_call
  ViaNow#run   !{ai, net}      <- now, so it inherited plan's effects
  ViaSend#fire !{}             <- send, so it did not
\end{term}

Remove the \eff{ai, net} declaration from \texttt{ViaNow.run} and it fails with
\texttt{declares \{\} but its body has \{ai, net\}}. \texttt{ViaSend.fire} may
stay \eff{}. The rule that \textbf{only the side that waits carries the
callee's effects} is visible right here.

\section{g7: \texttt{now} / \texttt{await} with deadlines}

\begin{lstlisting}
// g7: now / await with a deadline
//
//     Without a deadline, a now or await waits forever if no reply comes.
//     An actor is one thread, so while it waits it cannot process any
//     other message -- it turns to stone.
//
//       now t.m(a) timeout <ms> else <expr>
//       await f    timeout <ms> else <expr>
//
//     The else branch must have the same type as the body.
class Slow {
  method work(k) : int !{time} {
    wait(300);              // takes 300 ms
    reply(k * 2);
  }
  method fast(k) : int !{} {
    reply(k * 2);
  }
}

var s = new Slow();

// returns in time
var a = now s.fast(21) timeout 1000 else 0;
print("fast = " ++ a);

// times out -> the value of the else branch
var b = now s.work(21) timeout 50 else 0;
print("slow(timed out) = " ++ b);

// awaiting a future with a deadline
var f = future s.fast(5);
var c = await f timeout 1000 else 0;
print("await = " ++ c);
\end{lstlisting}

\begin{term}
fast = 42
slow(timed out) = 0
await = 10
\end{term}

\texttt{work} takes 300\,ms, so the \texttt{now} with a 50\,ms deadline times
out and becomes the \texttt{0} of the \texttt{else} branch. \texttt{fast} makes
the 1000\,ms deadline and returns \texttt{42}.

This sample also passes under \texttt{AIOS\_STRICT\_DEADLINE=1} (every wait has
a deadline). Change \texttt{else 0} to \texttt{else "zero"} and you get:

\begin{term}
[Type error] now: the else branch has type string but the call returns int
\end{term}

\section{g8: equality, unary minus, boolean literals}

\begin{lstlisting}
// g8: equality, unary minus, boolean literals
//
//     The lexer and the evaluator had all of these, but the expression
//     grammar did not, so they could not be written.
//
//       ==  !=          equality, on int / float / string / bool
//       - e             unary minus; lowers to neg(e)
//       true  false     boolean literals
//
//     Precedence is equality < comparison < ++ < additive < multiplicative
//     < unary minus, so 1 + 2 == 3 reads as (1 + 2) == 3.
class C {
  var flag = false;                     // a bool field
  method set(v) : unit !{mut} { flag = v; }
  method get() : bool !{} { reply(flag); }
  method t() : unit !{log, mut} {
    flag = true;
    if (flag == true)  { print("literal true ok"); }
    if (flag != false) { print("literal false ok"); }
    var b = 1 < 2;
    print("b = " ++ b);
    print("not-eq: " ++ (true != false));
  }
}
var c = new C();
send c.t();
\end{lstlisting}

\begin{term}
literal true ok
literal false ok
b = true
not-eq: true
\end{term}

All of these existed in the lexer and the evaluator but had \emph{no expression
grammar rule}, so they could not be written; the previous edition listed them
as a known gap. \texttt{flag} is initialized with \texttt{false}, so it is
\ty{bool}, and the parameter of \texttt{set(v)} is wired to \ty{bool} as well.

\section{g9: passing an actor as an argument}

\begin{lstlisting}
// g9: passing an actor as a message argument
//
//     With the parameter annotation `method m(x: D)` you can send to the
//     actor you received. Without it the parameter type is undetermined
//     and you get `send target is not actor`.
class D { method ping() : unit !{log} { print("pong"); } }
class P {
  method viaAnno(x: D) : unit !{log} {
    print("got actor");
    send x.ping();
    print("after send");
  }
}
var d = new D();
var p = new P();
send p.viaAnno(d);
\end{lstlisting}

\begin{term}
got actor
after send
pong
\end{term}

The parameter annotation \texttt{d: Device} fixes the class of that argument;
passing anything else is a type error. This annotation also adds an edge for
obligation levels (\S\ref{sec:levels}), which is exactly why \textbf{levels did
not need to be carried on actor types}.

\section{g10: two classes with a field of the same name}

\begin{lstlisting}
// g10: two different classes may have a field of the same name
//
//      A class body is checked under that class's own fields, so
//      Fork's held and Latch's held are different things.
class Fork {
  var held = 0;
  method take() : int !{mut} { held = 1; reply(held); }
}
class Latch {
  var held = 0;
  method arm()  : int !{mut} { held = 2; reply(held); }
}
var f = new Fork();
var l = new Latch();
print("fork  = " ++ (now f.take() timeout 200 else -1));
print("latch = " ++ (now l.arm()  timeout 200 else -1));
\end{lstlisting}

\begin{term}
fork  = 1
latch = 2
\end{term}

\begin{quote}\small
This is a regression test. The OCaml implementation used to \emph{stack} fields
into the type environment, so an assignment to the second class's \texttt{held}
was rejected as an overloaded name (Python and JavaScript were correct). Of the
586 programs at hand, 119 contained two classes with a field of the same name,
but in almost all of them another error came first and hid it.
\end{quote}

\section{Negative samples}

\texttt{docs/samples/reply\_inference/} holds samples whose purpose is to
confirm that they \emph{do not} pass.

\begin{center}
\small
\begin{tabular}{@{}ll@{}}
\toprule
File & What happens \\
\midrule
\texttt{s2\_missing\_reply}      & the result of a forgotten \texttt{reply} is used in arithmetic \\
\texttt{s3\_conflict}            & branches \texttt{reply} with different types \\
\texttt{s5\_overload\_ambiguity} & \texttt{a + b} with no principal type \\
\texttt{s8\_annotation\_conflict}& a \texttt{reply} that contradicts the declaration \\
\texttt{s9\_missing\_path}       & some path does not \texttt{reply} \\
\texttt{s10\_expose\_unannotated}& an exposed actor without annotations \\
\texttt{s12\_service\_exposed}   & the same, in the exposed integration sample \\
\texttt{s14\_select\_pattern}    & a \texttt{case} arity that differs from the signature \\
\texttt{s15\_double\_reply}      & a double \texttt{reply} \\
\midrule
\multicolumn{2}{@{}l@{}}{directly under \texttt{docs/samples/}} \\
\texttt{soundness2\_effect\_gap} & \texttt{future}+\texttt{await} escapes the effect check (a known hole) \\
\bottomrule
\end{tabular}
\end{center}

\part{Building and running}
\label{part:tool}

\section{Building the implementation}

\begin{term}
cd src && make thread_repl
\end{term}

Building with \texttt{dune} fails because \texttt{src/lexer.ml} is generated
twice, so use \texttt{src/Makefile}.

\section{Running a program}

Pass a file of REPL commands with \texttt{-f}.

\begin{term}
cat > run.repl <<'EOF'
load docs/samples/guide/g7_deadline.aipl
compile
EOF
./src/abclrepl_thread -q -f run.repl < /dev/null
\end{term}

The main REPL commands are
\texttt{load} / \texttt{compile} / \texttt{list} / \texttt{send} /
\texttt{ast} / \texttt{pprint} / \texttt{script} / \texttt{exit}.

\section{Running only the type checker}
\label{sec:tc}

The REPL depends on SDL, so a small driver exists for type checking alone.

\begin{term}
ocamlfind ocamlc -package unix -thread -linkpkg -w -a -I src -o tc \
  src/location.cmo src/types.cmo src/ast.cmo src/typing_env.cmo \
  src/infer.cmo src/typecheck.cmo src/parser.cmo src/lexer.cmo \
  docs/samples/reply_inference/tc_main.ml

./tc docs/samples/guide/g3_actors.aipl
\end{term}

It prints the table of inferred and declared return types, the effects of each
method, and the signature of each class.

\section{The type/run-time differential test}

There is a harness that compares the return type the checker assigned with the
type of the value actually \texttt{reply}ed at run time.

\begin{term}
python3 scripts/type_runtime_diff.py abclc/*.aipl docs/samples/guide/*.aipl
\end{term}

With \texttt{AIOS\_TYPE\_TRACE=1} the implementation prints
\texttt{[rtype] Class\#method = tag} on every \texttt{reply}. This mechanism
actually found an inconsistency where integer arithmetic returned a \ty{float}
($\vdash e : \ty{int}$ yet $e \Downarrow$ \texttt{VFloat}).
A deliberately mismatching negative sample can be marked as known with
\texttt{// TYPE\_RUNTIME\_DIFF: expect-mismatch}.

\textbf{Run it whenever you change the types.} It is the standing device for
catching a break in preservation.

\section{Environment variables}

\begin{center}
\begin{tabular}{@{}ll@{}}
\toprule
Variable & Effect \\
\midrule
\texttt{AIOS\_QUIET=1}            & suppress debug output such as type schemes \\
\texttt{AIOS\_STRICT\_DEADLINE=1} & make a \texttt{now} / \texttt{await} without a deadline an error \\
\texttt{AIOS\_LAX\_OVERLOAD=1}    & make an ambiguous overload a warning instead of an error \\
\texttt{AIOS\_LAX\_EXPOSE=1}      & make a missing annotation on an exposed actor a warning \\
\texttt{AIOS\_TYPE\_TRACE=1}      & print an \texttt{[rtype]} line on every \texttt{reply} \\
\texttt{AIOS\_MODEL\_PROVIDER}    & choose the provider for AI calls \\
\texttt{ABCL\_AI\_PROVIDER}       & the same (old name) \\
\texttt{REMOTE\_REVIEWER\_HOSTPORT} & the target of \texttt{remote\_review} \\
\bottomrule
\end{tabular}
\end{center}

\section{Running on real hardware --- Xinu / Raspberry Pi 3 and Pi 5}
\label{sec:xinu}

AIPL programs also run on a \textbf{bare-metal Raspberry Pi 3} (Embedded Xinu):
no operating system underneath, no libc, and the kernel itself carries the
actor machinery.

\subsection{Three runtimes, and what each is for}

\begin{center}
\small
\begin{tabular}{@{}lll@{}}
\toprule
Runtime & What it does & Checks \\
\midrule
OCaml REPL (canonical) & checks types, effects, deadlines, the seven disciplines & \textbf{yes} \\
Host VM (Mac) & runs \texttt{.avm}; values are readable & no \\
Xinu hardware (Pi 3) & the in-kernel VM runs \texttt{.avm} & no \\
Xinu hardware (Pi 5) & AIPL is lowered to C and JITed to AArch64 in-kernel & no \\
\bottomrule
\end{tabular}
\end{center}

This is the most important point in this section. \textbf{The \texttt{.avm} path
does not type-check.} Return-type annotations (\texttt{: string}), obligation
levels (\texttt{@3}) and effects (\texttt{!\{log\}}) are \textbf{accepted as
syntax and then discarded}. Every static guarantee comes from the canonical
(OCaml) side; the hardware is where \textbf{something already checked} is run.
So the workflow has two stages: pass the type check of \S\ref{sec:tc} first,
then lower to \texttt{.avm} and ship it to the board.

\subsection{Compiling to \texttt{.avm} and loading it}

An \texttt{.avm} file is integer actor bytecode (the first four bytes are
\texttt{AVM1}), and the format is \textbf{the same} on the host VM and on the
Pi 3.

\begin{term}
# 1. .aipl -> .avm
./_build/default/compile_avm.exe prog.aipl prog.avm

# 2. check the values on the host VM first
./_build/default/server.exe 8099 --no-open
curl --data-binary @prog.avm "http://127.0.0.1:8099/actor/loadvm?ask=0"
curl "http://127.0.0.1:8099/api/console"

# 3. send it to the board (the Pi 3 listens on 8080)
curl --data-binary @prog.avm "http://192.168.3.50:8080/actor/loadvm?ask=0"
loadvm: body=131 accepted=1 spawned actor id=1
\end{term}

\textbf{No kernel reflash is needed.} The \texttt{.avm} module is loaded into
the running kernel and starts as an actor there and then. The SD card is only
swapped when a \textbf{new VM instruction} has been added.

\begin{itemize}
\item \textbf{Without \texttt{?ask=0} nothing comes back.} Forget it and the
  request sits for 25 seconds and closes with zero bytes received. It looks
  exactly like a dead board, so suspect this first.
\item The board is \textbf{fragile under back-to-back HTTP}; leave 2--4 seconds
  between requests.
\end{itemize}

\subsection{Reading values back from the board}
\label{sec:pi3-console}

\texttt{print} on the Pi 3 goes to the serial line and the screen.
\textbf{Since 2026-09-04 the recent output can also be read with
\texttt{GET /api/console}}, in the same shape the host VM uses,
\texttt{\{"total":N,"lines":[...]\}}; add \texttt{?clear=1} to empty it after reading.

\begin{term}
curl --data-binary @g1_hello.avm "http://192.168.3.50:8080/actor/loadvm?ask=0"
loadvm: body=168 accepted=1 spawned actor id=1

curl "http://192.168.3.50:8080/api/console"
{"total":3,"lines":["a2: hello, AIPL","a2: tick 1","a2: tick 2"]}
\end{term}

\textbf{Until this went in, g5 was the only sample whose output had been checked
on the Pi 3 itself.} The Pi 5 returns the output in the \texttt{POST /cc} reply,
so the two boards could not be compared by running the same program on both.
On the first day it was readable, two defects surfaced that had been invisible
until then (a double \texttt{init} call, and actors from the previous program
staying alive --- \S\ref{sec:pi3-onemodule}).
\textbf{What you cannot observe, you cannot call correct.}

To read a return value only, external exposure still works (\S\ref{sec:expose}).

\lstinputlisting{samples/guide/g5_web.aipl}

\begin{term}
curl "http://192.168.3.50:8080/api/routes"
/echo -> a6

curl "http://192.168.3.50:8080/api/x/echo?method=say&args=hi"
echo: hi
\end{term}

\texttt{args} is URL-decoded before it is passed in (\texttt{\%20} and
\texttt{+} become spaces). The reply is awaited for up to 90 seconds, and the
wait ends as soon as the reply arrives.

\subsection{An on-board model --- \texttt{ai\_call}}
\label{sec:xinu-ai}

The Pi 3 kernel has \textbf{a small language model baked into it}.
\texttt{ai\_call} does not reach out to a service; it \textbf{runs inference on
the board}.

\begin{center}
\small
\begin{tabular}{@{}ll@{}}
\toprule
Model & \texttt{stories260K} (Llama 2 shape; \texttt{dim}=64, 5 layers, 8 heads, vocab 512) \\
Size & about 260K parameters, 1.06\,MB in fp32 \\
Where & the \texttt{.rodata} of the kernel image (kernel: 1.09\,MB $\to$ 2.16\,MB) \\
Arithmetic & soft float, D-cache off, 32-bit A53 \\
Measured & 2--6 seconds per \texttt{ai\_call} round trip (24 tokens generated) \\
\bottomrule
\end{tabular}
\end{center}

\begin{term}
curl "http://192.168.3.50:8080/api/x/sage?method=ask&args=Once%20upon%20a%20time"
MODEL[, there was a little girl named Lily. She loved to play outside in the p]END
\end{term}

The effect table (\S\ref{sec:builtins}) gives \texttt{ai\_call} the effect
\eff{ai, net}, yet on the Pi 3 \textbf{nothing leaves the board}. The effect
checker over-approximates, so the fact that \eff{net} does not actually occur
here does not compromise its soundness: it may forbid too much, never too
little.

\textbf{The host VM (Mac) has no model.} Calling \texttt{ai\_call} there
returns a string prefixed with \texttt{(no local model on host)}. It does not
manufacture a plausible-looking answer --- the absence of a model is written
into the value rather than quietly hidden.

\subsection{What the hardware path does not have}

The \texttt{.avm} path accepts a \textbf{subset} of the canonical language.
What is missing \textbf{fails on the spot}; nothing is dropped silently.

\begin{center}
\small
\begin{tabular}{@{}lll@{}}
\toprule
Feature & \texttt{.avm} path & Note \\
\midrule
Classes, fields, \texttt{send} / \texttt{send!} & yes & \\
Method-local \texttt{var}, \texttt{reply} & yes & locals become hidden fields \\
\texttt{if} / \texttt{while} & yes & \textbf{\texttt{while (c) do} --- parentheses required} \\
\texttt{now} / \texttt{future} / \texttt{await} / deadlines & yes & via \textbf{continuation splitting} (below) \\
\texttt{select} / \texttt{case} / \texttt{timeout} & yes & a \texttt{case} name must be a real method \\
Strings (values, \texttt{++}, arguments) & yes & \\
\texttt{web\_listen} / \texttt{web\_expose} & yes & \\
\texttt{ai\_call} & yes & on-board model (\S\ref{sec:xinu-ai}) \\
\midrule
Type / effect / level annotations & accepted, \textbf{discarded} & checked on the canonical side \\
\ty{float} and float literals & no & values are integers and strings only \\
Arrays (\texttt{array\_*}) & no & \texttt{avm: unsupported expression} \\
Math builtins (\texttt{sqrt} etc.) & no & likewise \\
\texttt{spawn} & no & likewise \\
\texttt{remote(...)} sends & no & \texttt{parse: unexpected token} \\
\texttt{replyto} & no & \texttt{avm: unknown variable} \\
\bottomrule
\end{tabular}
\end{center}

A word on \textbf{continuation splitting}. The VM has no instruction for a
synchronous call that returns a value. So wherever \texttt{now},
\texttt{await} or \texttt{select} appears, the method is \textbf{cut in two}
and the remainder is emitted as a separate method. To carry values across the
cut, method-local \texttt{var}s are turned into \textbf{hidden fields of the
actor}. The destination of \texttt{reply} needs the same care: after a split,
\texttt{sender} is no longer the original caller, so the caller is saved into
\texttt{\_\_snd} on entry to the method. The behaviour a user sees matches the
canonical implementation, but \textbf{more methods are generated} than were
written.

\subsection{The Pi 3 holds exactly one module at a time}
\label{sec:pi3-onemodule}

The Pi 3 VM has \textbf{a single module slot}
(\texttt{static unsigned char vm\_mod[...]}). \texttt{/actor/loadvm}
\textbf{overwrites it}, so \textbf{only the most recently loaded module is
live}.

This bit silently for a long time. Routes registered earlier stayed in the
table; only the actors behind them stopped working --- calling one got no reply,
and \texttt{vm\_web\_call} sat out its full 90 seconds and returned empty. After
loading thirteen modules back to back, the \texttt{/sage} and \texttt{/echo}
exposed first had both gone quiet, while an actor loaded at that moment
answered instantly.

\textbf{On 2026-09-04 this was brought in line with the Pi 5.} A \texttt{POST /cc}
replaces the Pi 5's whole actor world; \texttt{/actor/loadvm} now likewise
\textbf{folds away the previous module's actors} and empties the route table
first (resident C actors are left alone). Without that, the old actors stayed
alive as Xinu threads and \textbf{ate the process table} --- running the ten
guide samples in order, the ninth (g9) stopped starting its actors and HTTP
wedged shortly after (measured). It also removes the stale method-name pointers
that the mailboxes still held into the replaced \texttt{vm\_mod}.

\begin{itemize}
\item \textbf{Do not use \texttt{kill()} to fold them.} An actor thread calls
  \texttt{alloc\_obj} inside the \texttt{spawn} opcode and holds
  \texttt{objects\_mu} there; killing it mid-way stops every later spawn, taking
  \texttt{webactor} and HTTP down with it (also measured). \textbf{Mark them and
  wake them} instead, and let each leave its own loop.
\item Exposed routes \textbf{live and die with the program}: after a load,
  \texttt{/api/routes} lists only the new one.
\item \textbf{Try one module at a time} on the hardware; check it before
  loading the next.
\end{itemize}

\subsection{The ten canonical guide samples on hardware}

Every one of g1--g10 from Part \ref{part:sample} lowers to \texttt{.avm}, loads
onto the Pi 3, and starts as an actor.
\textbf{On 2026-09-04 all ten were checked on the hardware down to their output}
(\S\ref{sec:pi3-console}). Before that only \texttt{accepted=1} --- ``loaded and
started'' --- had been observed, plus the output of the exposed samples g5 and
\texttt{/sage}.

\begin{center}
\small
\begin{tabular}{@{}lll@{}}
\toprule
 & Canonical & Pi 3 hardware \\
\midrule
g1  & \texttt{hello, AIPL / tick 1 / tick 2} & matches \\
g2  & \texttt{twice = 43 / awaited = 20 / v=7} & matches \\
g3  & \texttt{ok, left=7 / sold out} & matches \\
g4  & \texttt{waiting / got 1 / via select:1} & matches \\
g5  & \texttt{/api/x/echo} $\to$ \texttt{echo: hi} & matches \\
g6  & \texttt{1 / 1 / n = 1} & matches \\
g7  & \texttt{fast = 42 / slow(timed out) = 0 / await = 10} & matches \\
g8  & \texttt{... / b = true / not-eq: true} & \textbf{\texttt{b = 1}} (below) \\
g9  & \texttt{got actor / after send / pong} & \textbf{order differs} (below) \\
g10 & \texttt{fork = 1 / latch = 2} & matches \\
\bottomrule
\end{tabular}
\end{center}

Neither mismatch is a gap in the front end; both come from the shape of the
execution vehicle.

\begin{itemize}
\item \textbf{g8's \texttt{b = 1}} --- this VM holds booleans as integers, and the
  value carries nothing that would print them back as \texttt{true} / \texttt{false}.
  \textbf{The Pi 5 prints \texttt{b = 1} too}, so this is not a difference between
  the boards.
\item \textbf{g9's ordering} --- Pi 3 actors are \textbf{real processes}, so the
  target of \texttt{send x.ping()} can run before \texttt{print("after send")}
  (the board printed \texttt{got actor / pong / after send}).
  The Pi 5 is a cooperative pump, so it always produces the canonical order.
  \textbf{Both are correct AIPL}: \texttt{send} promises no ordering.
\end{itemize}

That comparison holds \textbf{for one module loaded at a time}; it does
\textbf{not} mean the ten run side by side (\S\ref{sec:pi3-onemodule}).

\subsection{The Pi 5 --- the other board}
\label{sec:xinu-pi5}

The Pi 5 (\texttt{192.168.3.101}) \textbf{works differently}. Instead of
interpreting bytecode, it lowers AIPL to C on the board and the in-kernel C
compiler \texttt{cc} \textbf{JITs that to AArch64 machine code}.

\begin{center}
\small
\begin{tabular}{@{}lll@{}}
\toprule
 & Pi 3 & Pi 5 \\
\midrule
Execution & in-kernel VM interprets \texttt{.avm} & AIPL$\to$C$\to$in-kernel \texttt{cc} JIT \\
Front end & \texttt{compile\_avm} on the Mac & \textbf{on the board}: \texttt{abcl2c} \\
Submission & \texttt{POST /actor/loadvm?ask=0} (8080) & \texttt{POST /cc} (80) \\
Reading output & \texttt{/api/console} (\S\ref{sec:pi3-console}) & in the reply \\
Reading return values & \texttt{web\_expose} + \texttt{/api/x/} & same, or in the reply \\
Modules resident & one (\textbf{a new one replaces it}) & one (a new one replaces it) \\
Actor vehicle & \textbf{real Xinu processes} (truly concurrent) & cooperative pump (one thread) \\
\texttt{select} & continuation split + interception & \textbf{same} \\
\texttt{future} / \texttt{await} & continuation split (not concurrent) & runs at once (not concurrent) \\
\texttt{wait(ms)} & a VM opcode; works on every path & \textbf{works on every path} \\
Ten guide samples & \textbf{8 match exactly} & \textbf{all ten match} \\
Value kinds & integers and strings & int, string, float, list \\
\texttt{ai\_call} & opcode \texttt{0x52} & runtime \texttt{v\_ai\_call} \\
\texttt{ai\_call} measured & 2--6 s & \textbf{0.42 s} \\
Kernel & 1.09\,MB $\to$ 2.16\,MB & 0.94\,MB $\to$ 2.03\,MB \\
\bottomrule
\end{tabular}
\end{center}

\subsubsection*{Running a program}

Because the front end \textbf{lives on the board}, you can post AIPL source
directly. If the first word is \texttt{class} it goes through \texttt{abcl2c};
otherwise it is treated as C. \textbf{Whitespace and comments are skipped
before that test}, so canonical sources that open with \texttt{//} pass through
unchanged.

\begin{term}
curl --data-binary @g1_hello.aipl "http://192.168.3.101/cc"
hello, AIPL
tick 1
tick 2
=> 0
\end{term}

The trailing \texttt{=> 0} is \texttt{main}'s return value. From the shell,
\texttt{cc <file>} / \texttt{make <file>} handle \texttt{.abcl} / \texttt{.aipl}
transparently. The in-kernel model can also be poked directly with
\texttt{llm <prompt>} (output goes to the serial line).

\subsubsection*{What the program becomes}

What you write stays AIPL, but the generated C is readable. Classes become
\textbf{numbers}, objects are elements of \texttt{g\_obj[]}, methods become
functions named \texttt{m\_<Class>\_<method>}, and a method call becomes a
numeric \texttt{dispatch}.

\begin{term}
struct Obj { int cls; int f[1]; };   /* f[] holds the fields */
struct Obj g_obj[2048];
int v_stock;                          /* top-level vars become globals */
int m_Stock_take(int self, int a0, int a1, int a2, int a3) { ... }
int __new_init(int id, int meth, int a0, int a1, int a2, int a3);
int __dl_call(int to, int meth, int a0, int a1, int a2, int a3, int ms, int elsev);
int dispatch(int self, int meth, int a0, int a1, int a2, int a3);
int main() { v_stock = v_int(g_spawn(0)); ... }
\end{term}

\textbf{In this \texttt{cc}, \texttt{int} is 8 bytes} (stated in
\texttt{cc/ccpriv.h}). That matches the width of AIPL's tagged 64-bit values,
so a string (a pointer into the heap) survives being held in an \texttt{int}.
Note that \textbf{a function takes at most 8 parameters}; a ninth gives
\texttt{cc: too many parameters (max 8)}.

Values go through runtime helpers --- \texttt{v\_int} \texttt{v\_str}
\texttt{v\_floatlit} \texttt{v\_add} \texttt{v\_sub} \texttt{v\_mul}
\texttt{v\_div} \texttt{v\_lt} \texttt{v\_le} \texttt{v\_eq} \texttt{v\_ne}
\texttt{v\_and} \texttt{v\_or} \texttt{v\_not} \texttt{v\_truthy}
\texttt{v\_int\_of} \texttt{v\_print} \texttt{v\_ai\_call} \texttt{v\_list\_*}.
\texttt{v\_add} concatenates when either side is a string, so \texttt{++} is
mapped onto it.

\subsubsection*{The in-kernel model}

\texttt{ai\_call} lowers to \texttt{v\_ai\_call}, which renders the prompt to a
string, calls \texttt{llm\_run}, copies the result into the value heap and
returns it as a string. The model is \textbf{the same one as the Pi 3}
(\S\ref{sec:xinu-ai}): \texttt{ai\_call("Once upon a time")} agrees to the byte
(72 of them) across the Pi 3, the Pi 5 and the reference build on the Mac.
Only the speed differs: \textbf{the Pi 5 takes 0.42 seconds per call}. The
Pi 3 takes 2--6 seconds when freshly booted, and took 16.6 seconds once
thirteen modules had been loaded and actors had piled up (generation yields to
the scheduler per token, so more resident actors means slower --- though
\textbf{that causal link has not been isolated}).

\subsubsection*{What the on-board front end (\texttt{abcl2c}) accepts}

\begin{center}
\small
\begin{tabular}{@{}p{20mm}p{102mm}@{}}
\toprule
Accepted & classes, fields, methods, \texttt{var}, assignment,
       \texttt{if}/\texttt{else}, \texttt{while} (\texttt{do} optional),
       \texttt{return} \\
     & \texttt{new} (arguments are passed to \texttt{init} when there is one),
       \texttt{send}, \texttt{now} (returns a value) \\
     & \textbf{\texttt{reply(e)}} --- the same as \texttt{return e} here, since
       \texttt{now o.m()} takes \texttt{dispatch}'s return value. It must be
       \textbf{the last statement of its block}; otherwise it is refused,
       because a \texttt{return} would change the meaning \\
     & \textbf{annotations} \texttt{: T} / legacy \texttt{-> T} /
       \texttt{!\{log, mut\}} / \texttt{@ 2} / \texttt{x: D} /
       \texttt{var n: int} --- \textbf{accepted and discarded} (checking is the
       canonical side's job) \\
     & \textbf{\texttt{++}}, \textbf{\texttt{true} / \texttt{false}}
       (\ty{bool} is an \ty{int} here, so 1 / 0), unary minus \\
     & \textbf{postfix deadlines}:
       \texttt{now t.m(a) timeout <ms> else <expr>} (see below) \\
     & \textbf{\texttt{web\_listen} / \texttt{web\_expose}} (see below) \\
     & \textbf{\texttt{wait(ms)}} --- \textbf{only on the shell
       \texttt{cc}/\texttt{make} path} (see below) \\
     & \texttt{print} / \texttt{puts}, \texttt{ai\_call} \\
     & \textbf{\texttt{future} / \texttt{await}} --- but \textbf{not
       concurrent}; they run on the spot (\S\ref{sec:pi5-future}) \\
     & \textbf{\texttt{select} / \texttt{case} / \texttt{timeout}} --- but they
       \textbf{wait, on the Pi 3's continuation scheme} (\S\ref{sec:pi5-select}) \\
\midrule
Refused & \texttt{replyto}, \texttt{spawn}, \texttt{remote}, arrays
       (\texttt{array\_*}), the maths builtins (\texttt{sqrt} and friends) \\
     & field initialisers must be \textbf{integer or boolean literals}
       (they used to become a silent 0 when they were anything else) \\
\bottomrule
\end{tabular}
\end{center}

Refusals name what was refused:

\begin{term}
not supported on this device: select
unknown function: sqrt
new: unknown class: Zzz
reply must be the last statement of its block on this device
\end{term}

\noindent\textbf{Nothing quietly turns into something else.}

\subsubsection*{The ten canonical samples on the board}

\textbf{All ten of g1--g10 now run unchanged} when posted to \texttt{POST /cc}.
Previously none of them did --- every one died on its return-type annotation.
But \textbf{``runs'' and ``matches the canonical semantics'' are not the same
thing}.

\begin{center}
\small
\begin{tabular}{@{}lll@{}}
\toprule
Sample & Output (\texttt{POST /cc}) & Canonical? \\
\midrule
g1 & \texttt{hello, AIPL} / \texttt{tick 1} / \texttt{tick 2} & yes \\
g2 & \texttt{twice = 43} / \texttt{awaited = 20} / \texttt{v=7} & yes \\
g3 & \texttt{ok, left=7} / \texttt{sold out} & yes \\
g5 & exposed, then \texttt{echo: hi} from \texttt{/api/x/echo} & yes \\
g6 & \texttt{1} / \texttt{1} / \texttt{n = 1} & yes \\
g8 & \texttt{literal true ok} / \texttt{literal false ok} / \texttt{b = 1} & yes \\
g9 & \texttt{got actor} / \texttt{after send} / \texttt{pong} & yes \\
g10 & \texttt{fork  = 1} / \texttt{latch = 2} & yes \\
g7 & \texttt{fast = 42} / \texttt{slow(timed out) = 0} / \texttt{await = 10} & yes \\
g4 & \texttt{waiting} / \texttt{got 1} / \texttt{via select:1} & yes \\
\bottomrule
\end{tabular}
\end{center}

\textbf{On 2026-09-04 all ten came to match over \texttt{POST /cc} alone.} The last
two holdouts were g4 (\texttt{select} did not wait, \S\ref{sec:pi5-select}) and g7
(\texttt{wait} was unavailable under \texttt{/cc}, \S\ref{sec:pi5-budget}); both were
\textbf{properties of the path}, not gaps in the front end.
\textbf{There is no longer any need to pick a path.}

g8 prints \texttt{b = 1} where the canonical answer is \texttt{b = true}: this
board holds \ty{bool} as an \ty{int}. \textbf{The Pi 3 prints \texttt{b = 1} too},
so it is not a difference between the boards.

\subsubsection*{Deadlines are checked afterwards}
\label{sec:pi5-deadline}

\texttt{now t.m(a) timeout <ms> else <expr>} is accepted, but it
\textbf{does not interrupt the call}. Execution here is a cooperative pump and
\texttt{dispatch} runs to completion, so the deadline can only be used as an
\textbf{after-the-fact test of whether the call overran}.

\begin{term}
now s.quick(7) timeout 1000 else -99   ->   7     (within the deadline)
now s.heavy(7) timeout    1 else -99   -> -99     (overran -> the else branch)
\end{term}

\textbf{The value agrees with the canonical semantics}, but the call itself
still runs to the end. A program that relies on a deadline to \emph{stay
responsive} does not get that here.

\subsubsection*{Exposing actors --- \texttt{/api/x/}}

\texttt{web\_expose("/echo", "echo")} makes an actor reachable at\\
\texttt{GET /api/x/<path>?method=<name>\&args=<text>}.

\begin{term}
curl --data-binary @g5_web.aipl "http://192.168.3.101/cc"
curl "http://192.168.3.101/api/x/"
  exposed routes (web_expose):
    /api/x/echo  -> actor #0
curl "http://192.168.3.101/api/x/echo?method=say&args=hi"
  => echo: hi
\end{term}

\begin{itemize}
\item \textbf{The port argument is ignored.} The board's HTTP server already
  listens on 80, which satisfies what \texttt{web\_listen} is for (opening the
  boundary). The board records the number and serves the exposed routes on 80.
\item \textbf{Exposure lives and dies with the program.} Loading the next
  program clears the routes. The Pi 3 failure mode --- \textbf{stale routes
  that survive while their target dies} (\S\ref{sec:pi3-onemodule}) --- does
  not happen here.
\item The second argument is the name of a top-level \texttt{var}. The board
  identifies actors by number, so the name is resolved at translation time.
  An unknown name is refused by name.
\end{itemize}

\subsubsection*{The runaway budget, and \texttt{wait}}
\label{sec:pi5-budget}

JITed code is \textbf{cut off on a wall clock} in case it runs away, and
\textbf{the budget depends on the path}.

\begin{center}
\small
\begin{tabular}{@{}lll@{}}
\toprule
Path & Compute budget & \texttt{wait(ms)} \\
\midrule
\texttt{POST /cc} & \textbf{5 s} & \textbf{works} (10 s total) \\
shell \texttt{cc} / \texttt{make} & \textbf{10 s} & \textbf{works} (same) \\
\texttt{/actor/load} & 250 ms & refused \\
\bottomrule
\end{tabular}
\end{center}

\textbf{Since 2026-09-04, \texttt{wait} works under \texttt{/cc} too}, because both
now run in \textbf{their own process on their own stack}: a runaway can only blow
that stack and get cut off, never freeze the board.

\textbf{This was not done by sleeping in place.} HTTP handling can be entered from
the network process \emph{or} from the window-manager loop (\texttt{NULLPROC}), and
\texttt{proc\_sleep\_us} on \texttt{NULLPROC} context-switches \textbf{to itself} when
nothing else is runnable --- so it does not sleep at all. Instead \texttt{/cc} adopts
the shell's shape: spawn a runner and spin \texttt{proc\_resched()} until it is gone.

\textbf{Sleeping is not charged to the budget.} The budget is there to catch
\textbf{runaway computation}, and letting \texttt{wait(300)} eat it would defeat that.
The \textbf{total sleep} per run is capped at 10 s instead.

\begin{term}
curl --data-binary @g7_deadline.aipl "http://192.168.3.101/cc"
  fast = 42
  slow(timed out) = 0        <- wait(300) took effect. 0.40 s
  await = 10
\end{term}

\textbf{There is a trap here that half-kills the board.} The only thing advancing
the runner is \textbf{the caller spinning \texttt{proc\_resched()}}, and
\texttt{proc\_preempt()} never switches away from \texttt{NULLPROC}. So a runner that
is given up on \textbf{never runs again}: it keeps the static stack, and every later
\texttt{cc} answers ``the previous run has not finished'' (measured on hardware).
Three guards close it: cap the total sleep below the 15 s give-up, mark an
abandoned runner ``return at once'', and let the next call spin for 3 s to clear it.

\subsubsection*{\texttt{future} / \texttt{await} are not concurrent}
\label{sec:pi5-future}

\texttt{future o.m(x)} \textbf{runs on the spot}: it calls \texttt{dispatch}
immediately, stores the value in a slot and returns a handle. \texttt{await h}
reads that slot. A deadline on \texttt{await} \textbf{can never fire}, because
the value is already there. In other words, \textbf{there is no ``fire it off
and collect it later''}.

The Pi 3 is the same in this respect --- continuation splitting does not run
anything concurrently either --- so \textbf{the two boards have the same
temperament here}.

\textbf{A genuinely concurrent implementation was written and then removed,
because it froze the board.} It gave each call its own Xinu process and had
\texttt{await} wait for it. The value heaps (\texttt{vheap\_alloc},
\texttt{lheap\_alloc}) were protected by masking interrupts, and that was not
enough: the runaway state \texttt{g\_deadline} / \texttt{g\_aborted}, the output
capture \texttt{g\_cap} and the single FIFO behind \texttt{cc\_pump()} were all
still shared. The safeguards in place (a wall-clock cap, a liveness check
before reusing a slot) \textbf{never fired} --- something lower down stops
first. \textbf{``Protect the allocator and you can go concurrent'' was wrong.}

\subsubsection*{\texttt{select} --- rebuilt on the Pi 3's continuation scheme}
\label{sec:pi5-select}

\textbf{On 2026-09-03--04 this was rebuilt to match the Pi 3.}
Until then the Pi 5's \texttt{select} \textbf{did not wait}: there is no way to
wait in a cooperative pump, so it merely scanned the messages already queued for
the actor. In g4, \texttt{job} has not arrived when \texttt{send w.serve()} runs
\texttt{serve}, so it \textbf{always fell into the timeout branch} --- the gap
this chapter used to call ``the one place where the Pi 3 is ahead''.

The Pi 3's scheme \textbf{waits without waiting}. It touches neither the
scheduler nor the execution vehicle.

\begin{itemize}
\item The method containing the \texttt{select} sets a hidden field
  \texttt{\_\_sel} to that \texttt{select}'s number and \textbf{ends there}.
\item \texttt{case job(k) -> \{ ... \}} is spliced in \textbf{at the top of method
  \texttt{job}}: if \texttt{\_\_sel} holds its number, clear the mark and run the
  \texttt{case} body followed by the rest of the \texttt{select}'s method;
  otherwise run \texttt{job}'s own body.
\item A \texttt{case}'s parameter maps \textbf{by position} onto the receiving
  method's parameter.
\end{itemize}

The message a \texttt{select} waits for therefore \textbf{arrives as an ordinary
method call}. No suspension, no mailbox scan.

\textbf{The \texttt{timeout} branch is decided by ``nobody came'', not by time.}
This device has no vehicle that can wait, so the branch runs once the FIFO has
been drained and no one claimed the \texttt{select} (\texttt{\_\_sel\_sweep});
the millisecond value is unusable. The Pi 3 has a separate \texttt{\_\_Timer}
actor do \texttt{wait(ms)}, so \textbf{there it really is time} --- but the
meaning, ``run the timeout branch if the awaited message never comes'', is the same.

One more change: \textbf{the FIFO is now drained right after a top-level
\texttt{send}}. Pi 3 actors are real processes, so the target runs first; the
Pi 5 only queued the message, and a following \texttt{now} entered the target's
method ahead of it, defeating the interception.

\begin{center}
\small
\begin{tabular}{@{}lll@{}}
\toprule
 & Pi 3 & Pi 5 \\
\midrule
How & continuation split + interception & \textbf{same} \\
Can it wait? & yes & \textbf{yes} \\
\texttt{timeout} decided by & a \texttt{wait(ms)} in another actor & the FIFO draining with no claim \\
g4's output & canonical & \textbf{canonical} \\
\bottomrule
\end{tabular}
\end{center}

\begin{term}
curl --data-binary @g4_select.aipl "http://192.168.3.101/cc"
waiting
got 1
via select:1            <- used to be "timed out"

# when nobody sends the awaited message
waiting
timed out
\end{term}

\subsubsection*{Bad input does not stop the board}

The front end has been swept with \textbf{mutation fuzzing} (ASan + UBSan,
4{,}000 cases). On the board, the eight inputs that used to freeze the kernel
(\texttt{future} / \texttt{select} / \texttt{replyto} / \texttt{ai\_cal} (a
typo) / an argument-carrying \texttt{new} on a class with no \texttt{init} /
\texttt{new} of an unknown class / a statement after \texttt{reply} / a
deadline) were all replayed: \textbf{every one returns a named error at once,
and the board survives every time}.

\texttt{new <unknown class>} is worth singling out. \texttt{class\_idx()}
returns $-1$ for a name it does not know, and that $-1$ used to go straight
into the syntax tree, so the emitter read \texttt{CLS[-1]} --- \textbf{an
out-of-bounds read on real hardware}. The fuzzer found it.


\part{Appendix}
\label{part:appendix}

\section{Grammar summary}

\begin{term}
program    ::= decl+

decl       ::= "class" ID "{" field* method+ "}"
             | "var" ID "=" expr ";"
             | "var" ID "=" "new" ID "(" args ")" ";"
             | "send"  target "." ID "(" args ")" ";"
             | "send!" target "." ID "(" args ")" ";"
             | ID "(" args ")" ";"

field      ::= "var"   ID "=" expr ";"
             | "float" ID "=" expr ";"

method     ::= "method" ID "(" params ")" ret? lvl? eff? "{" stmt+ "}"
params     ::= (param ("," param)*)?
param      ::= ID | ID ":" ClassName | ID ":" "reply"
ret        ::= ":" ("int"|"float"|"string"|"bool"|"unit"|"any"|ClassName)
lvl        ::= "@" INT                    -- obligation level; inferred if omitted
eff        ::= "!" "{" (ID ("," ID)*)? "}"

target     ::= ID
             | "remote" "(" STRING "," STRING ")"

stmt       ::= ID "=" expr ";"
             | "var" ID "=" expr ";"
             | "var" ID "=" "new" ID "(" args ")" ";"
             | ID "(" args ")" ";"
             | "call" ID "(" args ")" ";"
             | "send"  ("self"|"sender"|target) "." ID "(" args ")" ";"
             | "send!" target "." ID "(" args ")" ";"
             | "if" "(" expr ")" stmt ("else" stmt)?
             | "while" expr "do" stmt
             | "{" stmt* "}"
             | "select" "{" case* timeout_case? "}"

case         ::= "case" ID "(" ids? ")" "->" "{" stmt+ "}"
timeout_case ::= "timeout" INT "->" "{" stmt+ "}"

expr       ::= INT | FLOAT | STRING | "true" | "false" | ID
             | "-" expr
             | expr ("=="|"!="|">"|"<"|">="|"<="|"++"|"+"|"-"|"*"|"/") expr
             | "new" ID "(" args ")"
             | "now" target "." ID "(" args ")" deadline?
             | "future" target "." ID "(" args ")"
             | "await" expr deadline?
             | "replyto"                  -- the reply destination as a value
             | ID "(" args ")"
             | "(" expr ")"

-- with else it is tau; without else it is result<tau>
deadline   ::= "timeout" INT "else" expr
             | "timeout" INT
\end{term}

Operator precedence, weakest first:
\texttt{==} \texttt{!=} $\;<\;$
\texttt{>} \texttt{<} \texttt{>=} \texttt{<=} $\;<\;$
\texttt{++} $\;<\;$
\texttt{+} \texttt{-} $\;<\;$
\texttt{*} \texttt{/} $\;<\;$ unary minus. Every binary operator is
\textbf{left associative} (equality included) --- \texttt{1 == 2 == false}
reads as \texttt{((1 == 2) == false)} and evaluates to \texttt{true}. Unary
minus is a prefix and can be stacked (\texttt{- - 3} is \texttt{3}).

\section{Environment variables of the implementation}
\label{sec:envvars}

Some variables change how checking and execution behave. Set them to
\texttt{1} to enable; the default is always ``not enabled''.

\begin{center}
\small
\begin{tabular}{@{}lp{9.4cm}@{}}
\toprule
Variable & Effect \\
\midrule
\texttt{AIOS\_STRICT\_DEADLINE} &
  make a \texttt{now} / \texttt{await} without a deadline an \textbf{error};
  by default it is only a warning. Set it if you want to stay inside the
  fragment for which the soundness theorem holds (AIPL$^{-2}$) \\
\texttt{AIOS\_LAX\_OVERLOAD} &
  restore the \textbf{old behaviour} for ambiguous overloads: warn and pick the
  default candidate. The default is strict (an error) \\
\texttt{AIOS\_LAX\_EXPOSE} &
  downgrade a missing annotation on an exposed actor (\S\ref{sec:expose}) to a
  warning. The default is an error \\
\texttt{AIOS\_TYPE\_TRACE} &
  print \texttt{[rtype] Class\#method = tag} on every \texttt{reply}; used by
  the differential test (\texttt{scripts/type\_runtime\_diff.py}) that compares
  the static return type with the type actually returned \\
\texttt{AIOS\_QUIET} &
  suppress informational messages from the checker and the REPL \\
\texttt{AIOS\_STRICT\_PROTOCOL} &
  warn about the ``unfinished'' part when the steps of a session cannot be laid
  out statically. Silent by default (\S\ref{sec:session}) \\
\texttt{AIOS\_STRICT\_RESOURCE} &
  report \texttt{acquire} calls whose name is not a string literal; the order
  cannot be checked there (\S\ref{sec:resorder}) \\
\texttt{AIOS\_SHOW\_LEVELS} &
  print the inferred obligation levels and the global resource order \\
\texttt{AIOS\_LAX\_WAIT} &
  downgrade the detection of circular waits from an error to a warning.
  The default is an error \\
\texttt{AIOS\_LAX\_ACTOR} &
  do not distinguish actor types (unify \texttt{actor(A)} with
  \texttt{actor(B)}). The default is to distinguish them \\
\texttt{AIOS\_LAX\_ASSIGN} &
  allow assignment to an undeclared name (the old behaviour).
  The default is an error \\
\texttt{AIOS\_MODEL\_PROVIDER} &
  the model vendor used by \texttt{ai\_call}; the default is \texttt{gemini} \\
\bottomrule
\end{tabular}
\end{center}

\noindent For instance, to work under the discipline of always writing deadlines:

\begin{lstlisting}
AIOS_STRICT_DEADLINE=1 ./src/abclrepl_thread -q -f run.repl
\end{lstlisting}

\section{Common pitfalls}

\begin{itemize}
\item \textbf{\texttt{now} is an expression.} \texttt{now x.m();} is a syntax
  error; always receive it, as in
  \texttt{var r = now x.m() timeout 100 else 0;}.
\item \textbf{The target of \texttt{send} is a variable}, not a class name.
\item \textbf{\texttt{now self.m()} cannot be written} (it would be a
  self-deadlock, so the grammar forbids it). To factor out a computation, move
  it to another actor.
\item \textbf{You cannot assign an actor to \texttt{var x = 0;} later}
  (it has been inferred as \ty{int}). Create it as
  \texttt{var x = new Foo();} from the start.
\item \textbf{String concatenation is \texttt{++}}; \texttt{+} is numeric only.
\item \textbf{Fields come before methods.}
\item \textbf{Write a deadline on every wait.} Otherwise you get a warning, and
  under \texttt{AIOS\_STRICT\_DEADLINE=1} it does not pass at all.
\item \textbf{Add \texttt{?ask=0} when loading onto hardware.} Without it
  \texttt{/actor/loadvm} never answers, which looks like a dead board
  (\S\ref{sec:xinu}).
\item \textbf{\texttt{print} on the board goes to the serial line.} To read a
  value over the network, \texttt{web\_expose} it and call \texttt{/api/x/}.
\item \textbf{Effects are checked once you declare them} (gradual). Declare
  nothing and nothing is said; declare something and the difference from the
  body is pointed at.
\end{itemize}

\section{Present limits}

\begin{itemize}
\item \textbf{What runs on the hardware (Xinu / Pi 3) is a subset of the
  canonical language.} No \ty{float}, arrays, math builtins, \texttt{spawn},
  \texttt{remote} or \texttt{replyto}; type, effect and level annotations are
  accepted and discarded (checking happens on the canonical side).
  See \S\ref{sec:xinu}.
\item Splitting a wait into \texttt{future} + \texttt{await} escapes effect
  propagation (the ``known hole'' in \S\ref{sec:effects}). Fixing it requires
  carrying effects in the \ty{future} type.
\item There is no variable declaration with a type annotation,
  \texttt{var x : T = e;}.
\item \texttt{sender} is \ty{any}, so \texttt{send sender.m(...)} checks
  neither the existence of the method nor its types.
\item The result of a remote send (\texttt{remote(...)}) stays \ty{any}. There
  is no mechanism for reading the interface description of the peer node.
\item Method parameters are monomorphic; polymorphic methods cannot be written.
\item There is no implicit conversion between \ty{int} and \ty{float}, and
  \texttt{a + b} with both operands unbound needs an annotation.
\item There are no array literals \texttt{[...]}; build arrays with
  \texttt{array\_empty} and \texttt{array\_push}.
\item A reply destination (\texttt{replyto}) \textbf{cannot be kept in a
  field}, though it can be passed as an argument (\S\ref{sec:reply1st}).
  A field's type comes from its initializer, so \texttt{var w = 0;} is \ty{int}
  and \texttt{w = replyto} is a type mismatch.
\item The global resource order can only be followed for \texttt{acquire} calls
  whose name is a \textbf{string literal} (\S\ref{sec:resorder}).
\item The static session check reaches only as far as \textbf{synchronous}
  callees (\S\ref{sec:session}).
\item Termination is not guaranteed. Deadlock freedom means ``nothing gets
  stuck'', not ``everything finishes''. \texttt{while 1 > 0 do \{\}} does not
  stop, but it is not a wait.
\end{itemize}

\section{Changes from the previous edition}
\label{sec:migrate-en}

\subsection*{2026-09-03 --- the Pi 5 accepts much more}

The language did not change. What changed is \textbf{what the Pi 5's on-board
front end (\texttt{abcl2c}) accepts}; see \S\ref{sec:xinu-pi5}.

\begin{center}
\small
\begin{tabular}{@{}lll@{}}
\toprule
 & Previous & This edition \\
\midrule
Canonical samples (unchanged, \texttt{POST /cc}) & \textbf{0 / 10} & \textbf{10 / 10 run} \\
\quad of which canonical & 0 & \textbf{10} \\
Annotations \texttt{: T} / \texttt{!\{...\}} / \texttt{@N} & rejected & accepted, discarded \\
\texttt{reply} & refused & same as \texttt{return} (block-final) \\
\texttt{++} / \texttt{true} / \texttt{false} / unary minus & absent & accepted \\
Postfix \texttt{timeout … else …} & refused & accepted (\textbf{checked after}) \\
\texttt{web\_listen} / \texttt{web\_expose} & refused & accepted (\texttt{/api/x/}) \\
\texttt{wait(ms)} & refused & accepted (\texttt{/cc} too, \S\ref{sec:pi5-budget}) \\
\texttt{future} / \texttt{await} & refused & accepted (\textbf{not concurrent}) \\
\texttt{select} / \texttt{case} & refused & accepted (\textbf{waits}) \\
\bottomrule
\end{tabular}
\end{center}

Several \textbf{silent breakages} were fixed along the way: \texttt{new} of an
unknown class read \texttt{CLS[-1]} (an out-of-bounds read on hardware);
top-level \texttt{var}s referenced from a method produced C that would not
compile; field initialisers silently became 0 when they were not numeric; and
string literals were truncated at 62 bytes, desynchronising the lexer.

\label{sec:migrate}

What changed since the previous edition (2026-07-30).

\subsection{Removed from the language}

\begin{center}
\small
\begin{tabular}{@{}lp{8.6cm}@{}}
\toprule
 & Content \\
\midrule
\texttt{become} & \textbf{removed from the language}. ABCL/1 does not have it;
  once restricted enough to keep types sound it becomes almost the same as
  writing a state field and a branch; and not one of 543 real programs used it.
  Write behaviour replacement with a state field and \texttt{if} \\
\bottomrule
\end{tabular}
\end{center}

\subsection{Added}

None of these break existing programs. Every annotation may be omitted, and
omitting it means it is inferred.

\begin{center}
\small
\begin{tabular}{@{}lp{7.4cm}l@{}}
\toprule
 & Content & Section \\
\midrule
Failure type & a wait with no \texttt{else} returns \texttt{result}$\langle\tau\rangle$
  & \S\ref{sec:result} \\
First-class reply & \texttt{replyto} / \texttt{answer} / parameter type \texttt{reply}, treated linearly
  & \S\ref{sec:reply1st} \\
Resource order & the \texttt{acquire} / \texttt{release} pairs, and the global order read from their nesting
  & \S\ref{sec:resorder} \\
Obligation levels & \texttt{@$n$}, inferred if omitted; also enforced across node borders
  & \S\ref{sec:levels} \\
Session types & the type checker reads the run-time protocol declarations
  & \S\ref{sec:session} \\
The \texttt{select} discipline & a mandatory deadline, and a check that a sender exists
  & \S\ref{sec:selectrule} \\
Mesh deployment & ship the source and type-check it at the destination before instantiating
  & \S\ref{sec:mesh} \\
\bottomrule
\end{tabular}
\end{center}

\subsection{Made stricter}

\begin{center}
\small
\begin{tabular}{@{}lp{8.6cm}@{}}
\toprule
 & Content \\
\midrule
Assignment to undeclared names & now \textbf{rejected}. Previously a new
  variable was silently created, so a misspelled field name produced no error
  at all (\texttt{AIOS\_LAX\_ASSIGN=1} restores the old behaviour) \\
Actor types & \texttt{actor(A)} and \texttt{actor(B)} are now distinguished
  (\texttt{AIOS\_LAX\_ACTOR=1} restores the old behaviour) \\
\texttt{reply} totality & a method waited on by \texttt{now}/\texttt{await}
  must \texttt{reply} on every path, whatever its return type \\
\bottomrule
\end{tabular}
\end{center}

\subsection{The file extension}

The source extension is \texttt{.aipl} (\texttt{.abcl} has been retired
completely). The implementation does not look at the extension, so old files
still load, but write new ones as \texttt{.aipl}.

\subsection{Core and extensions --- the language has two layers}

There are three implementations, but they are \textbf{not the same language}.

\begin{center}
\small
\begin{tabular}{@{}lp{4.0cm}p{5.4cm}@{}}
\toprule
Layer & Implemented by & Content \\
\midrule
\textbf{Core} & all three & the specification this guide describes \\
\textbf{Extensions} & the Python implementation only & record types, row
  polymorphism, refinement types via \texttt{where}, CSP channels, ownership
  (\texttt{pub}/\texttt{move}), tuples, \texttt{match}, saga,
  \texttt{scope \{ future ... \}} \\
\bottomrule
\end{tabular}
\end{center}

\textbf{This guide specifies the core only.} Which layer a program belongs to
can be determined mechanically with \texttt{tools/classify\_corpus.sh}
(see \texttt{docs/AIPL\_LAYERS.md}).

\section{Related documents}

\begin{itemize}
\item \texttt{docs/aipl\_paper5\_ja.pdf} ---
  \textbf{the paper (5th edition, latest)}. The current specification, and what
  was taken from (or left out of) ABCL/1 and Kobayashi's type systems, with the
  reasons and the trade-offs.
\item \texttt{coq/AIPLSoundness2.v} ---
  \textbf{the machine-checked proof (2nd version)}, in Coq/Rocq. For the core
  \emph{including} \texttt{await}, type soundness, type safety and deadlock
  freedom are proved with no axioms. The header also discusses the
  \textbf{largest syntax} for which all three hold at once.
\item \texttt{coq/AIPLSoundness.v} --- the 1st version (deadlock freedom only
  for the fragment without \texttt{await}).
\item \texttt{docs/AIPL\_LAYERS.md} ---
  \textbf{AIPL has two layers}: the definitions of core (all three
  implementations) and extensions (Python only), and the classification by
  \texttt{tools/classify\_corpus.sh}.
\item \texttt{docs/aipl\_reply\_inference\_ja.pdf} --- return type inference
  from \texttt{reply}, the rationale for annotations, and measurements.
\item \texttt{docs/aipl\_vs\_abcl1\_ja.pdf} --- a comparison with ABCL/1.
\item \texttt{docs/aipl\_abcl1\_features\_ja.pdf} --- the features inherited
  from ABCL/1.
\item \texttt{docs/aipl\_kobayashi\_synthesis\_ja.pdf} --- a design proposal
  for taking in the results of Kobayashi's work.
\item \texttt{docs/ABCLCP\_TYPE\_SOUNDNESS\_REPORT.md} --- type soundness
  (Markdown).
\item \texttt{docs/AIOS\_LANGUAGE\_DESIGN.md}, \texttt{docs/CAPABILITIES.md} ---
  the design as AIOS and the list of capabilities.
\item \texttt{docs/samples/reply\_inference/README.md} --- what the negative
  samples are for.
\end{itemize}

\end{document}
